\section{AMC 12 12A 2026}
        \begin{enumerate}
        	\item (\textit{AMC 12 12A 2026}) A sign in a parking lot states that the total price includes a $27\%$ tax. To the nearest whole number, what percentage of the amount a customer pays is this tax?


 $\textbf{(A)}\   21\%  \qquad 	\textbf{(B)}\   22\%  \qquad 	\textbf{(C)}\   27\%  \qquad 	\textbf{(D)}\   37\%  \qquad 	\textbf{(E)}\   73\%$

 \textbf{Answer: }A



\textbf{Solution 1}

Let the original ticket be $\$100$. Then the tax would be $\$100 \cdot 27\% = \$27$. The customer pays $100+27=\$127$, so the percentage is $\frac{27}{127}=\fbox{\textbf{(A)}\ 21\%}$.

 



\textbf{Solution 2}

Let $a$ represent the cost of the ticket. Then, the total cost (with tax) is $\frac{127}{100} a$. The tax is $\frac{27}{100}a$. Thus tax takes up $\frac{27}{127}$. Through long division to three decimals, this is 0.213, or 21.3\%, and so the answer is $\fbox{\textbf{(A)}\ 21\%}$.Note: in this problem, don't take the shortcut by estimating unless you're really good at it.

 



\textbf{Solution 3 (estimation of 27/127)}

Let us say that you want to estimate $\frac{27}{127}$ because high risk, high reward (so basically you're a maniac).

 We wish to figure out what percent of 127 is 27. First, we know that $10\%$ of $127$ is $12.7$. Multiplying by two gives $20\%$ of $127$, which we find to be $25.4$. This means we are very close to the answer. Because $1\%$ of $127$ is $1.27$, we know that $21\%$ is $26.67$. On the contrary, $22\%$ is $1.27$ more, so it would be farther away from $27$ than $26.67$. Thus, our answer is $\fbox{\textbf{(A)}\ 21\%}$.

 Note: the problem where this is linked is covered 1:00-1:20

 


	\item (\textit{AMC 12 12A 2026}) Square $ABCD$ has side length $24$. Let $E$ be the midpoint of $\overline{BC}$, and let $F$ be the intersection of $\overline{AE}$ and $\overline{BD}$.


 
                    \begin{figure}[H]
                        \centering
                        \includegraphics[width=0.5\textwidth]{Problem/AMC_12_12A_2026_2_Problem_Figure_1.png}
                    \end{figure}
                    


 What is the area of $\triangle BEF$?


 $\textbf{(A)}\   36  \qquad 	\textbf{(B)}\   40  \qquad 	\textbf{(C)}\   45  \qquad 	\textbf{(D)}\   48  \qquad 	\textbf{(E)}\   52$

 \textbf{Answer: }D



\textbf{Solution 1}

To find the intersection point of the two lines, we can turn them into a system of equations. $\overline{AFE}$ becomes $y=-2x+24$, and $\overline{BFD}$ becomes $y=x$. Solving, we get $y=x=8$. Therefore the height of the triangle is $8$ and the base is $12$. So $\frac{8\cdot12}{2}=\boxed{48}$.~MrThinker

 
Edit: Both $\overline{AFE}$ and $\overline{BFD}$ were written as $\overline{ABC}$ ~ GETTHEMONKEY

 



\textbf{Solution 2}

$\overline{AD}$ is parallel to $\overline{BC}$, so $\angle{ADF} = \angle{FBE}$ and $\angle{DAF} = \angle{FEB}$. By the AA similarity theorem, $ADF \sim EBF$. Let the height of the triangle $EBF$ be $x$. $AD = 24$ and$EB = 12$, so the height of the triangle $DAF$ is twice the height of $FEB$. The sum of the two heights covers the whole height of the rectangle, so $3x = 24$, and $x = 8$. Therefore, the area of the triangle is $\frac{12 \cdot8}{2} = 48$ $\fbox{D}$. $\newline$$\sim$cyrus825

 



\textbf{Solution 3}

We can see that $\overline{AE}$ is a median of $\triangle ABC$. Let $\overline{BD}$ intersect $\overline{AC}$ at $G$, so $\overline{BG}$ is another median and $F$ is the centroid. The medians of $\triangle ABC$ split it into six triangles of equal area, and $\triangle BEF$ is one of them. Therefore,
 \[[BEF]=\frac{1}{6}[ABC]=\frac{1}{6}\cdot288=\boxed{\textbf{(D)}~48}\]




\textbf{Solution 4 (Coordinate geometry)}

We can solve this problem using coordinate geometry. Define B to be $(0,0)$. We can find the intersection of lines $\overline{BD}$ and $\overline{AE}$ by defining these lines as $y=x$ and $y=-2x+24$, respectively. Solving these equations, we find the intersection point to be at $(8,8)$ Finally, using Shoelace Theorem, we find the total area to be $\frac{1}{2}|(0+0+12\cdot 8)-(0+0+0)|=48$, so the answer is $\boxed{\textbf{D}}$~Voidling

 



\textbf{Solution 5 (Cheese)}

Measure the base and height, and scale it to match the quantities in the problem. Then calculate the area. WARNING!!! - Only do this if you have 30 seconds left.$\newline$$\sim$MATH\_1S\_AWESOME

 Edit: Diagrams are NOT drawn to scale so this won't work$\newline$Edit: Diagrams are usually drawn to scale, even though they say it doesn't have to be drawn to scale so you don't use the niche technique of ruler bashing$\newline$Edit: why would you have 30 seconds left while on the second problem

 



\textbf{Solution 6 (Area ratios)}

We know that triangle $ABE$ has area $\frac{12 \times 24}{2}$, which is just $144$. Since the line $BD$ bisects the angle $\angle ABE$, we can use the properties of the angle bisector theorem. By the angle bisector theorem:
 \[\frac{AF}{FE} = \frac{AB}{BE}\]
 \[\frac{AF}{FE} = \frac{24}{12} = \frac{2}{1}\]Therefore, $AF = 2FE$. This means that $FE = \frac{1}{3}AE$, which helps us conclude that the area of $\triangle BFE$ is $\frac{1}{3}$ of the area of $\triangle ABE$. 

 So our answer is:
 \[\frac{144}{3} = \boxed{48}\]




\textbf{Solution 7 (Bashing)}

We know that square $ABCD$ has area of 576, triangle $ABE$ has area of 144 and triangle $BCD$ has area 288. Now we use the answer choices to figure out the area of the remaining two triangles. Using 48 as the area of triangle $BEF$ we can find that the triangle’s height to be 8. This means triangle $AFD$ has a height of 24 - 8 = 16. Now we calculate the area of triangle $AFD$ to be 192. Since triangle $ABE$ and $BCD$ overlap we must subtract the area of $BEF$. We get 144 + 288 + 192 - 48 =576. Since 576 is the area of $ABCD$ we know the area of $BEF$ is 48. Hence we get answer 
 \[\boxed{48}\]




\textbf{Solution 8 (super duper cheese)}

Recognize this is the \href{https://tasks.illustrativemathematics.org/content-standards/tasks/1571}{construction of a third.}

 $\frac{24}{3} \times 12 \times \frac{1}{2} = \boxed{\textbf{(D)}~48}$




\textbf{Solution 9 (Cross Product)}

Establish a coordinate system with point $B$ at the origin $(0,0)$, $C$ at $(24,0)$, $A$ at $(0,24)$, and $D$ at $(24,24)$. Since $E$ is the midpoint of $\overline{BC}$, its coordinates are $(12,0)$.

 This allows us to write linear equations representing lines $AE$ and $BD$:

 
            \begin{itemize}
        \item Line $AE$ passes through $(0,24)$ and $(12,0)$, giving $y = -2x + 24$.
        \item Line $BD$ passes through $(0,0)$ and $(24,24)$, giving $y = x$.
            \end{itemize}
            Setting these equations equal to find their intersection point $F$:
 \[x = -2x + 24 \implies 3x = 24 \implies x = 8\]Thus, the intersection $F$ is at $(8,8)$. Since point $B$ is at $(0,0)$, we can represent two sides of $\triangle BEF$ as vectors in three dimensions (assuming a $z$-component of $0$):
 \[\overrightarrow{BF} = \langle 8, 8, 0 \rangle \quad \text{and} \quad \overrightarrow{BE} = \langle 12, 0, 0 \rangle\]

 Taking the cross product of these two vectors gives:
 \[\overrightarrow{BF} \times \overrightarrow{BE} = \begin{vmatrix} \mathbf{i} & \mathbf{j} & \mathbf{k} \\ 8 & 8 & 0 \\ 12 & 0 & 0 \end{vmatrix} = \langle 0, 0, (8)(0) - (8)(12) \rangle = \langle 0, 0, -96 \rangle\]

 Taking half the magnitude of the cross product gives the area of $\triangle BEF$:
 \[\text{Area} = \frac{1}{2} \left\| \overrightarrow{BF} \times \overrightarrow{BE} \right\| = \frac{1}{2} \sqrt{0^2 + 0^2 + (-96)^2} = \frac{96}{2} = \boxed{\textbf{(D) } 48}\]~thesuperioraxolotl

 


	\item (\textit{AMC 12 12A 2026}) Half of the internal angles of a polygon measure $150^{\circ}$ each, and the other half of its internal angles measure $190^{\circ}$ each. How many sides does the polygon have?


 $\textbf{(A)}\   36  \qquad 	\textbf{(B)}\   40  \qquad 	\textbf{(C)}\   60  \qquad 	\textbf{(D)}\   72  \qquad 	\textbf{(E)}\   180$

 \textbf{Answer: }A



\textbf{Solution 1}

Let the polygon have n sides. The sum of the interior angles of this polygon is $(n-2)\cdot 180$. $\frac{n}{2}$ of the angles are $150$ and $\frac{n}{2}$ of the angles are $190$, so
 \[\frac{n}{2} \cdot 150 + \frac{n}{2} \cdot 190= (n-2)\cdot 180\]
 \[75n + 95n = 180n - 360\]
 \[10n = 360\]
 \[n = 36\]Therefore, our answer is $\boxed{\textbf{(A)}\hspace{1mm}36}$$\newline$$\sim$cyrus825

 Slight $\LaTeX$ corrections by ~singularity1

 



\textbf{Solution 2}

Because we know that half of the angles are $190^c$ and the other half is $150^c$, the average angle is $\frac{190^c+150^c}{2} = 170^c$ on average for each angle. If the average is $170^c$, then 
 \[\frac{180*(n-2)}{n} = 170\]
 \[170n=180*(n-2) \to 0 = 10n-360\]
 \[10n=360 \to n = \boxed{36}\]




\textbf{Solution 3}

$\frac{lcm(340,360)}{360}=34 \Rightarrow 34+2= \boxed{\textbf{(A)}\hspace{1mm}36}$

 The calculation is trivial using prime factors

 



\textbf{Solution 4}

We can find the average exterior angle by doing $180^\circ-150^\circ=30^\circ$ for the first angle and $180^\circ-190^\circ=-10^\circ$ for the second. Since there are an equal amount of both exterior angles, we can find the average exterior angle:
 \[\frac{30^\circ+\left(-10^\circ\right)}{2} = 10^\circ\]Since the sum of the exterior angles of any polygon is $360^\circ$
 \[\frac{360^\circ}{10^\circ}=\boxed{\textbf{(A)}\hspace{1mm}36}\]




\textbf{Solution 4}

Let the polygon have $2n$ sides. Then $n$ angles measure $150^\circ$ and $n$ angles measure $190^\circ$ meaning that the total interior angle sum is $150n+190n=340n$ Since every polygon has interior angle sum $180(Sides-2)$, the polygon has interior angle sum $180(2n-2)=360n-360$. Therefore, $340n=360n-360$, so $20n=360$ meaning $n=18$ Hence, the polygon has $2n=\fbox{36}$ sides.

 


	\item (\textit{AMC 12 12A 2026}) Let $a_{1}, a_{2}, a_{3}$, $\dots$ be a sequence with $a_{1} = 94$, $a_{2} = 190$, and $a_{n}$ equal to the arithmetic mean of the previous $n-1$ terms for $n \ge 3$. What is the value of $a_{10}$?


 $\textbf{(A)}\   47  \qquad 	\textbf{(B)}\   95  \qquad 	\textbf{(C)}\   142  \qquad 	\textbf{(D)}\   284  \qquad 	\textbf{(E)}\   564$

 \textbf{Answer: }C



\textbf{Solution 1}

By the definition of the sequence, $a_3 = \frac{94+190}{2} = 142$, $a_4 = \frac{94+190+142}{3} = 142$, and so on. Trying a few more terms keeps giving us 142 and it never changes, so the answer is $\boxed{\textbf{(C)}\hspace{1mm}142}$$\newline$$\sim$cyrus825

 minor $\LaTeX$ edits by ~singularity1

 



\textbf{Solution 2}

For $n>3,$ we have 
 \[a_{n-1} = \dfrac{a_{n-2}+a_{n-3}+...+a_1}{n-2}\implies a_{n-2}+a_{n-3}+...+a_1=(n-2)a_{n-1}.\]

 Then, 
 \[a_{n} = \dfrac{a_{n-1}+a_{n-2}+...+a_1}{n-1}=\dfrac{a_{n-1}+(n-2)a_{n-1}}{n-1}=a_{n-1}.\]

 Since $a_3 = \frac{94+190}{2}=142,$ we have 
 \[a_{10}=a_9=a_8=...=a_3=\boxed{\textbf{(C) }142}\]~Tacos\_are\_yummy\_1

 



\textbf{Solution 3 (Analyzing Option Choices)}

One can see that the average of $94$ and $190$ will be between $94$ and $190$, so $a_10$ will obviously be between $94$ and $190$. Looking at the option choices, we can see that only $95$ and $142$ are in between $94$ and $190$, and the average is in between $94$ and $190$, so cannot be $95$. Thus, the answer must be $\boxed{142}$




\textbf{Solution 4 (Understanding Arithmetic Mean)}

It is known that adding a value equal to the arithmetic mean into a data set does not change the overall mean. A quick proof of the previous statement is as follows:

 Let the mean of a set of $n$ numbers be $\overline{x}$. Adding $\overline{x}$ to the set of $n$ numbers yields the new sum: $n \cdot \overline{x} \rightarrow (n+1)\cdot \overline{x}$.

 See that the total number of values increases from $n$ to $n+1$, and thus the new mean of the set of $n+1$ values gives $\frac{(n+1)\cdot \overline{x}}{n+1}= \overline{x}$.

 Hence, in our problem, we need only find the arithmetic mean of the two initial values of the sequence. That is: $\frac{a_{1} + a_{2}}{2} = \frac{94+190}{2} = \boxed{142}$. 

 


 


	\item (\textit{AMC 12 12A 2026}) How many ordered triples of positive integers $(x,y,z)$ satisfy both of the following equations?
<mjx-container class="MathJax CtxtMenu\_Attached\_0" ctxtmenu\_counter="0" display="true" jax="CHTML" style="font-size: 117.4\%; position: relative;" tabindex="0"><mjx-math aria-hidden="true" class="MJX-TEX" display="true" style="margin-left: 0px; margin-right: 0px;"><mjx-mtable style="min-width: 7.305em;"><mjx-table><mjx-itable><mjx-mtr><mjx-mtd style="text-align: right; padding-right: 0px; padding-bottom: 0.15em;"><mjx-mi class="mjx-i"><mjx-c class="mjx-c1D465 TEX-I"></mjx-c></mjx-mi><mjx-mo class="mjx-n" space="3"><mjx-c class="mjx-c2B"></mjx-c></mjx-mo><mjx-mi class="mjx-i" space="3"><mjx-c class="mjx-c1D466 TEX-I"></mjx-c></mjx-mi><mjx-mo class="mjx-n" space="3"><mjx-c class="mjx-c2B"></mjx-c></mjx-mo><mjx-mi class="mjx-i" space="3"><mjx-c class="mjx-c1D467 TEX-I"></mjx-c></mjx-mi><mjx-tstrut></mjx-tstrut></mjx-mtd><mjx-mtd style="text-align: left; padding-left: 0px; padding-bottom: 0.15em;"><mjx-mi class="mjx-n"></mjx-mi><mjx-mo class="mjx-n" space="4"><mjx-c class="mjx-c3D"></mjx-c></mjx-mo><mjx-mn class="mjx-n" space="4"><mjx-c class="mjx-c33"></mjx-c><mjx-c class="mjx-c31"></mjx-c></mjx-mn><mjx-tstrut></mjx-tstrut></mjx-mtd></mjx-mtr><mjx-mtr><mjx-mtd style="text-align: right; padding-right: 0px; padding-top: 0.15em;"><mjx-mi class="mjx-i"><mjx-c class="mjx-c1D465 TEX-I"></mjx-c></mjx-mi><mjx-mo class="mjx-n" space="3"><mjx-c class="mjx-c2B"></mjx-c></mjx-mo><mjx-mn class="mjx-n" space="3"><mjx-c class="mjx-c33"></mjx-c></mjx-mn><mjx-mi class="mjx-i"><mjx-c class="mjx-c1D466 TEX-I"></mjx-c></mjx-mi><mjx-mo class="mjx-n" space="3"><mjx-c class="mjx-c2B"></mjx-c></mjx-mo><mjx-mn class="mjx-n" space="3"><mjx-c class="mjx-c33"></mjx-c></mjx-mn><mjx-mi class="mjx-i"><mjx-c class="mjx-c1D467 TEX-I"></mjx-c></mjx-mi><mjx-tstrut></mjx-tstrut></mjx-mtd><mjx-mtd style="text-align: left; padding-left: 0px; padding-top: 0.15em;"><mjx-mi class="mjx-n"></mjx-mi><mjx-mo class="mjx-n" space="4"><mjx-c class="mjx-c3D"></mjx-c></mjx-mo><mjx-mn class="mjx-n" space="4"><mjx-c class="mjx-c34"></mjx-c><mjx-c class="mjx-c31"></mjx-c></mjx-mn><mjx-tstrut></mjx-tstrut></mjx-mtd></mjx-mtr></mjx-itable></mjx-table></mjx-mtable></mjx-math><mjx-assistive-mml display="block" unselectable="on"><math display="block" xmlns="http://www.w3.org/1998/Math/MathML"><mtable columnalign="right left" columnspacing="0em" displaystyle="true" rowspacing="3pt"><mtr><mtd><mi>x</mi><mo>+</mo><mi>y</mi><mo>+</mo><mi>z</mi></mtd><mtd><mi></mi><mo>=</mo><mn>31</mn></mtd></mtr><mtr><mtd><mi>x</mi><mo>+</mo><mn>3</mn><mi>y</mi><mo>+</mo><mn>3</mn><mi>z</mi></mtd><mtd><mi></mi><mo>=</mo><mn>41</mn></mtd></mtr></mtable></math></mjx-assistive-mml></mjx-container>
$\textbf{(A)}\   4  \qquad 	\textbf{(B)}\   5  \qquad 	\textbf{(C)}\   6  \qquad 	\textbf{(D)}\   7  \qquad 	\textbf{(E)}\   8$

 \textbf{Answer: }A



\textbf{Solution 1}

Subtracting the first equation from the second equation gives
 \[2y+2z=10\]
 \[y+z=5\]By the first equation, $x = 26$. The ordered pairs of positive integers $(y, z)$ such that $y + z = 5$ are $(1, 4)$, $(2, 3)$, $(3, 2)$, and $(4, 1)$. Therefore, there are $4$ ordered triples, or $\boxed{\textbf{(A)}}$$\newline$$\sim$cyrus825, minor $\LaTeX$ fixes and polishing by Voidling

 



\textbf{Solution 1.5}

After getting $y+z=5$, we use stars and bars to see that there are $6$ non-negative options for $y$ and $z$. After subtracting the 2 cases where y or z equals $0$, our final answer is $6-2=4$ —> $\boxed{\textbf{(A)}}$

 -Amon26

 



\textbf{Solution 2}

Let $S = y + z$. Substituting into both equations gives the system $x + S = 31$ and $x + 3S = 41$.Equating expressions for $x$ yields $31 - S = 41 - 3S \implies 2S = 10 \implies S = 5$, which fixes $x = 26$.Since $y, z \in \mathbb{Z}^+$ must satisfy $y + z = 5$, we apply the Stars and Bars theorem for $k = 2$ parts.The number of positive integer solutions $(y, z)$ is $\binom{5 - 1}{2 - 1} = \binom{4}{1} = 4$.Because $x = 26$ is uniquely determined, there are exactly $\boxed{4}$ ordered triples $(x, y, z)$.

 



\textbf{Solution 3 (Gaussian Elimination/RREF)}

The given system of two equations in three unkowns can be represented as an augmented $3 \times 4$ coefficient matrix. We apply Gaussian Elimination (row reduction) to produce an equivalent matrix in Reduced Row Echelon Form (RREF).

 We begin by representing our system as the augmented $3 \times 4$ matrix:

 \begin{bmatrix}{1} & {1} & {1} & {31} \\{1} & {3} & {3} & {41}\end{bmatrix}

 Apply Gaussian Elimination yields the resulting RREF matrix:

 \begin{bmatrix}{1} & {0} & {0} & {26} \\{0} & {1} & {1} & {5}\end{bmatrix}

 Hence, we have that $x=26$ and that $y+z=5$. The only ordered pairs $(y,z)$ for which $y+z=5$ are $(1,4)$, $(2,3)$, $(3,2)$, and $(4,1)$. Thus, there are $\boxed{4}$ such ordered triples. 

 



\textbf{Solution 0 (Brute Force the Brute Force)}

Note: If you are planning to read this solution, it will probably hurt your brain.

 There are \(\binom{30}{2}=435\) possible ordered triples of positive integers \((x,y,z)\) that sum to 31. A mere 435 possibilities suggests the use of brute force to calculate and check all 435 outcomes in the 75 - minute exam.$\newline$(1,1,29): 91 != 41 $\newline$(1,2,28): 91 != 41 $\newline$(1,3,27): 91 != 41 $\newline$(1,4,26): 91 != 41 $\newline$(1,5,25): 91 != 41 $\newline$(1,6,24): 91 != 41 $\newline$(1,7,23): 91 != 41 $\newline$(1,8,22): 91 != 41 $\newline$(1,9,21): 91 != 41 $\newline$(1,10,20): 91 != 41 $\newline$(1,11,19): 91 != 41 $\newline$(1,12,18): 91 != 41 $\newline$(1,13,17): 91 != 41 $\newline$(1,14,16): 91 != 41 $\newline$(1,15,15): 91 != 41 $\newline$(1,16,14): 91 != 41 $\newline$(1,17,13): 91 != 41 $\newline$(1,18,12): 91 != 41 $\newline$(1,19,11): 91 != 41 $\newline$(1,20,10): 91 != 41 $\newline$(1,21,9): 91 != 41 $\newline$(1,22,8): 91 != 41 $\newline$(1,23,7): 91 != 41 $\newline$(1,24,6): 91 != 41 $\newline$(1,25,5): 91 != 41 $\newline$(1,26,4): 91 != 41 $\newline$(1,27,3): 91 != 41 $\newline$(1,28,2): 91 != 41 $\newline$(1,29,1): 91 != 41 $\newline$

 (2,1,28): 89 != 41 $\newline$(2,2,27): 89 != 41 $\newline$(2,3,26): 89 != 41 $\newline$(2,4,25): 89 != 41 $\newline$(2,5,24): 89 != 41 $\newline$(2,6,23): 89 != 41 $\newline$(2,7,22): 89 != 41 $\newline$(2,8,21): 89 != 41 $\newline$(2,9,20): 89 != 41 $\newline$(2,10,19): 89 != 41 $\newline$(2,11,18): 89 != 41 $\newline$(2,12,17): 89 != 41 $\newline$(2,13,16): 89 != 41 $\newline$(2,14,15): 89 != 41 $\newline$(2,15,14): 89 != 41 $\newline$(2,16,13): 89 != 41 $\newline$(2,17,12): 89 != 41 $\newline$(2,18,11): 89 != 41 $\newline$(2,19,10): 89 != 41 $\newline$(2,20,9): 89 != 41 $\newline$(2,21,8): 89 != 41 $\newline$(2,22,7): 89 != 41 $\newline$(2,23,6): 89 != 41 $\newline$(2,24,5): 89 != 41 $\newline$(2,25,4): 89 != 41 $\newline$(2,26,3): 89 != 41 $\newline$(2,27,2): 89 != 41 $\newline$(2,28,1): 89 != 41 $\newline$

 (3,1,27): 87 != 41 $\newline$(3,2,26): 87 != 41 $\newline$(3,3,25): 87 != 41 $\newline$(3,4,24): 87 != 41 $\newline$(3,5,23): 87 != 41 $\newline$(3,6,22): 87 != 41 $\newline$(3,7,21): 87 != 41 $\newline$(3,8,20): 87 != 41 $\newline$(3,9,19): 87 != 41 $\newline$(3,10,18): 87 != 41 $\newline$(3,11,17): 87 != 41 $\newline$(3,12,16): 87 != 41 $\newline$(3,13,15): 87 != 41 $\newline$(3,14,14): 87 != 41 $\newline$(3,15,13): 87 != 41 $\newline$(3,16,12): 87 != 41 $\newline$(3,17,11): 87 != 41 $\newline$(3,18,10): 87 != 41 $\newline$(3,19,9): 87 != 41 $\newline$(3,20,8): 87 != 41 $\newline$(3,21,7): 87 != 41 $\newline$(3,22,6): 87 != 41 $\newline$(3,23,5): 87 != 41 $\newline$(3,24,4): 87 != 41 $\newline$(3,25,3): 87 != 41 $\newline$(3,26,2): 87 != 41 $\newline$(3,27,1): 87 != 41 $\newline$

 (4,1,26): 85 != 41 $\newline$(4,2,25): 85 != 41 $\newline$(4,3,24): 85 != 41 $\newline$(4,4,23): 85 != 41 $\newline$(4,5,22): 85 != 41 $\newline$(4,6,21): 85 != 41 $\newline$(4,7,20): 85 != 41 $\newline$(4,8,19): 85 != 41 $\newline$(4,9,18): 85 != 41 $\newline$(4,10,17): 85 != 41 $\newline$(4,11,16): 85 != 41 $\newline$(4,12,15): 85 != 41 $\newline$(4,13,14): 85 != 41 $\newline$(4,14,13): 85 != 41 $\newline$(4,15,12): 85 != 41 $\newline$(4,16,11): 85 != 41 $\newline$(4,17,10): 85 != 41 $\newline$(4,18,9): 85 != 41 $\newline$(4,19,8): 85 != 41 $\newline$(4,20,7): 85 != 41 $\newline$(4,21,6): 85 != 41 $\newline$(4,22,5): 85 != 41 $\newline$(4,23,4): 85 != 41 $\newline$(4,24,3): 85 != 41 $\newline$(4,25,2): 85 != 41 $\newline$(4,26,1): 85 != 41 $\newline$

 (5,1,25): 83 != 41 $\newline$(5,2,24): 83 != 41 $\newline$(5,3,23): 83 != 41 $\newline$(5,4,22): 83 != 41 $\newline$(5,5,21): 83 != 41 $\newline$(5,6,20): 83 != 41 $\newline$(5,7,19): 83 != 41 $\newline$(5,8,18): 83 != 41 $\newline$(5,9,17): 83 != 41 $\newline$(5,10,16): 83 != 41 $\newline$(5,11,15): 83 != 41 $\newline$(5,12,14): 83 != 41 $\newline$(5,13,13): 83 != 41 $\newline$(5,14,12): 83 != 41 $\newline$(5,15,11): 83 != 41 $\newline$(5,16,10): 83 != 41 $\newline$(5,17,9): 83 != 41 $\newline$(5,18,8): 83 != 41 $\newline$(5,19,7): 83 != 41 $\newline$(5,20,6): 83 != 41 $\newline$(5,21,5): 83 != 41 $\newline$(5,22,4): 83 != 41 $\newline$(5,23,3): 83 != 41 $\newline$(5,24,2): 83 != 41 $\newline$(5,25,1): 83 != 41 $\newline$

 (6,1,24): 81 != 41 $\newline$(6,2,23): 81 != 41 $\newline$(6,3,22): 81 != 41 $\newline$(6,4,21): 81 != 41 $\newline$(6,5,20): 81 != 41 $\newline$(6,6,19): 81 != 41 $\newline$(6,7,18): 81 != 41 $\newline$(6,8,17): 81 != 41 $\newline$(6,9,16): 81 != 41 $\newline$(6,10,15): 81 != 41 $\newline$(6,11,14): 81 != 41 $\newline$(6,12,13): 81 != 41 $\newline$(6,13,12): 81 != 41 $\newline$(6,14,11): 81 != 41 $\newline$(6,15,10): 81 != 41 $\newline$(6,16,9): 81 != 41 $\newline$(6,17,8): 81 != 41 $\newline$(6,18,7): 81 != 41 $\newline$(6,19,6): 81 != 41 $\newline$(6,20,5): 81 != 41 $\newline$(6,21,4): 81 != 41 $\newline$(6,22,3): 81 != 41 $\newline$(6,23,2): 81 != 41 $\newline$(6,24,1): 81 != 41 $\newline$

 (7,1,23): 79 != 41 $\newline$(7,2,22): 79 != 41 $\newline$(7,3,21): 79 != 41 $\newline$(7,4,20): 79 != 41 $\newline$(7,5,19): 79 != 41 $\newline$(7,6,18): 79 != 41 $\newline$(7,7,17): 79 != 41 $\newline$(7,8,16): 79 != 41 $\newline$(7,9,15): 79 != 41 $\newline$(7,10,14): 79 != 41 $\newline$(7,11,13): 79 != 41 $\newline$(7,12,12): 79 != 41 $\newline$(7,13,11): 79 != 41 $\newline$(7,14,10): 79 != 41 $\newline$(7,15,9): 79 != 41 $\newline$(7,16,8): 79 != 41 $\newline$(7,17,7): 79 != 41 $\newline$(7,18,6): 79 != 41 $\newline$(7,19,5): 79 != 41 $\newline$(7,20,4): 79 != 41 $\newline$(7,21,3): 79 != 41 $\newline$(7,22,2): 79 != 41 $\newline$(7,23,1): 79 != 41 $\newline$

 (8,1,22): 77 != 41 $\newline$(8,2,21): 77 != 41 $\newline$(8,3,20): 77 != 41 $\newline$(8,4,19): 77 != 41 $\newline$(8,5,18): 77 != 41 $\newline$(8,6,17): 77 != 41 $\newline$(8,7,16): 77 != 41 $\newline$(8,8,15): 77 != 41 $\newline$(8,9,14): 77 != 41 $\newline$(8,10,13): 77 != 41 $\newline$(8,11,12): 77 != 41 $\newline$(8,12,11): 77 != 41 $\newline$(8,13,10): 77 != 41 $\newline$(8,14,9): 77 != 41 $\newline$(8,15,8): 77 != 41 $\newline$(8,16,7): 77 != 41 $\newline$(8,17,6): 77 != 41 $\newline$(8,18,5): 77 != 41 $\newline$(8,19,4): 77 != 41 $\newline$(8,20,3): 77 != 41 $\newline$(8,21,2): 77 != 41 $\newline$(8,22,1): 77 != 41 $\newline$

 (9,1,21): 75 != 41 $\newline$(9,2,20): 75 != 41 $\newline$(9,3,19): 75 != 41 $\newline$(9,4,18): 75 != 41 $\newline$(9,5,17): 75 != 41 $\newline$(9,6,16): 75 != 41 $\newline$(9,7,15): 75 != 41 $\newline$(9,8,14): 75 != 41 $\newline$(9,9,13): 75 != 41 $\newline$(9,10,12): 75 != 41 $\newline$(9,11,11): 75 != 41 $\newline$(9,12,10): 75 != 41 $\newline$(9,13,9): 75 != 41 $\newline$(9,14,8): 75 != 41 $\newline$(9,15,7): 75 != 41 $\newline$(9,16,6): 75 != 41 $\newline$(9,17,5): 75 != 41 $\newline$(9,18,4): 75 != 41 $\newline$(9,19,3): 75 != 41 $\newline$(9,20,2): 75 != 41 $\newline$(9,21,1): 75 != 41 $\newline$

 (10,1,20): 73 != 41 $\newline$(10,2,19): 73 != 41 $\newline$(10,3,18): 73 != 41 $\newline$(10,4,17): 73 != 41 $\newline$(10,5,16): 73 != 41 $\newline$(10,6,15): 73 != 41 $\newline$(10,7,14): 73 != 41 $\newline$(10,8,13): 73 != 41 $\newline$(10,9,12): 73 != 41 $\newline$(10,10,11): 73 != 41 $\newline$(10,11,10): 73 != 41 $\newline$(10,12,9): 73 != 41 $\newline$(10,13,8): 73 != 41 $\newline$(10,14,7): 73 != 41 $\newline$(10,15,6): 73 != 41 $\newline$(10,16,5): 73 != 41 $\newline$(10,17,4): 73 != 41 $\newline$(10,18,3): 73 != 41 $\newline$(10,19,2): 73 != 41 $\newline$(10,20,1): 73 != 41 $\newline$

 (11,1,19): 71 != 41 $\newline$(11,2,18): 71 != 41 $\newline$(11,3,17): 71 != 41 $\newline$(11,4,16): 71 != 41 $\newline$(11,5,15): 71 != 41 $\newline$(11,6,14): 71 != 41 $\newline$(11,7,13): 71 != 41 $\newline$(11,8,12): 71 != 41 $\newline$(11,9,11): 71 != 41 $\newline$(11,10,10): 71 != 41 $\newline$(11,11,9): 71 != 41 $\newline$(11,12,8): 71 != 41 $\newline$(11,13,7): 71 != 41 $\newline$(11,14,6): 71 != 41 $\newline$(11,15,5): 71 != 41 $\newline$(11,16,4): 71 != 41 $\newline$(11,17,3): 71 != 41 $\newline$(11,18,2): 71 != 41 $\newline$(11,19,1): 71 != 41 $\newline$

 (12,1,18): 69 != 41 $\newline$(12,2,17): 69 != 41 $\newline$(12,3,16): 69 != 41 $\newline$(12,4,15): 69 != 41 $\newline$(12,5,14): 69 != 41 $\newline$(12,6,13): 69 != 41 $\newline$(12,7,12): 69 != 41 $\newline$(12,8,11): 69 != 41 $\newline$(12,9,10): 69 != 41 $\newline$(12,10,9): 69 != 41 $\newline$(12,11,8): 69 != 41 $\newline$(12,12,7): 69 != 41 $\newline$(12,13,6): 69 != 41 $\newline$(12,14,5): 69 != 41 $\newline$(12,15,4): 69 != 41 $\newline$(12,16,3): 69 != 41 $\newline$(12,17,2): 69 != 41 $\newline$(12,18,1): 69 != 41 $\newline$

 (13,1,17): 67 != 41 $\newline$(13,2,16): 67 != 41 $\newline$(13,3,15): 67 != 41 $\newline$(13,4,14): 67 != 41 $\newline$(13,5,13): 67 != 41 $\newline$(13,6,12): 67 != 41 $\newline$(13,7,11): 67 != 41 $\newline$(13,8,10): 67 != 41 $\newline$(13,9,9): 67 != 41 $\newline$(13,10,8): 67 != 41 $\newline$(13,11,7): 67 != 41 $\newline$(13,12,6): 67 != 41 $\newline$(13,13,5): 67 != 41 $\newline$(13,14,4): 67 != 41 $\newline$(13,15,3): 67 != 41 $\newline$(13,16,2): 67 != 41 $\newline$(13,17,1): 67 != 41 $\newline$

 (14,1,16): 65 != 41 $\newline$(14,2,15): 65 != 41 $\newline$(14,3,14): 65 != 41 $\newline$(14,4,13): 65 != 41 $\newline$(14,5,12): 65 != 41 $\newline$(14,6,11): 65 != 41 $\newline$(14,7,10): 65 != 41 $\newline$(14,8,9): 65 != 41 $\newline$(14,9,8): 65 != 41 $\newline$(14,10,7): 65 != 41 $\newline$(14,11,6): 65 != 41 $\newline$(14,12,5): 65 != 41 $\newline$(14,13,4): 65 != 41 $\newline$(14,14,3): 65 != 41 $\newline$(14,15,2): 65 != 41 $\newline$(14,16,1): 65 != 41 $\newline$

 (15,1,15): 63 != 41 $\newline$(15,2,14): 63 != 41 $\newline$(15,3,13): 63 != 41 $\newline$(15,4,12): 63 != 41 $\newline$(15,5,11): 63 != 41 $\newline$(15,6,10): 63 != 41 $\newline$(15,7,9): 63 != 41 $\newline$(15,8,8): 63 != 41 $\newline$(15,9,7): 63 != 41 $\newline$(15,10,6): 63 != 41 $\newline$(15,11,5): 63 != 41 $\newline$(15,12,4): 63 != 41 $\newline$(15,13,3): 63 != 41 $\newline$(15,14,2): 63 != 41 $\newline$(15,15,1): 63 != 41 $\newline$

 (16,1,14): 61 != 41 $\newline$(16,2,13): 61 != 41 $\newline$(16,3,12): 61 != 41 $\newline$(16,4,11): 61 != 41 $\newline$(16,5,10): 61 != 41 $\newline$(16,6,9): 61 != 41 $\newline$(16,7,8): 61 != 41 $\newline$(16,8,7): 61 != 41 $\newline$(16,9,6): 61 != 41 $\newline$(16,10,5): 61 != 41 $\newline$(16,11,4): 61 != 41 $\newline$(16,12,3): 61 != 41 $\newline$(16,13,2): 61 != 41 $\newline$(16,14,1): 61 != 41 $\newline$

 (17,1,13): 59 != 41 $\newline$(17,2,12): 59 != 41 $\newline$(17,3,11): 59 != 41 $\newline$(17,4,10): 59 != 41 $\newline$(17,5,9): 59 != 41 $\newline$(17,6,8): 59 != 41 $\newline$(17,7,7): 59 != 41 $\newline$(17,8,6): 59 != 41 $\newline$(17,9,5): 59 != 41 $\newline$(17,10,4): 59 != 41 $\newline$(17,11,3): 59 != 41 $\newline$(17,12,2): 59 != 41 $\newline$(17,13,1): 59 != 41 $\newline$

 (18,1,12): 57 != 41 $\newline$(18,2,11): 57 != 41 $\newline$(18,3,10): 57 != 41 $\newline$(18,4,9): 57 != 41 $\newline$(18,5,8): 57 != 41 $\newline$(18,6,7): 57 != 41 $\newline$(18,7,6): 57 != 41 $\newline$(18,8,5): 57 != 41 $\newline$(18,9,4): 57 != 41 $\newline$(18,10,3): 57 != 41 $\newline$(18,11,2): 57 != 41 $\newline$(18,12,1): 57 != 41 $\newline$

 (19,1,11): 55 != 41 $\newline$(19,2,10): 55 != 41 $\newline$(19,3,9): 55 != 41 $\newline$(19,4,8): 55 != 41 $\newline$(19,5,7): 55 != 41 $\newline$(19,6,6): 55 != 41 $\newline$(19,7,5): 55 != 41 $\newline$(19,8,4): 55 != 41 $\newline$(19,9,3): 55 != 41 $\newline$(19,10,2): 55 != 41 $\newline$(19,11,1): 55 != 41 $\newline$

 (20,1,10): 53 != 41 $\newline$(20,2,9): 53 != 41 $\newline$(20,3,8): 53 != 41 $\newline$(20,4,7): 53 != 41 $\newline$(20,5,6): 53 != 41 $\newline$(20,6,5): 53 != 41 $\newline$(20,7,4): 53 != 41 $\newline$(20,8,3): 53 != 41 $\newline$(20,9,2): 53 != 41 $\newline$(20,10,1): 53 != 41 $\newline$

 (21,1,9): 51 != 41 $\newline$(21,2,8): 51 != 41 $\newline$(21,3,7): 51 != 41 $\newline$(21,4,6): 51 != 41 $\newline$(21,5,5): 51 != 41 $\newline$(21,6,4): 51 != 41 $\newline$(21,7,3): 51 != 41 $\newline$(21,8,2): 51 != 41 $\newline$(21,9,1): 51 != 41 $\newline$

 (22,1,8): 49 != 41 $\newline$(22,2,7): 49 != 41 $\newline$(22,3,6): 49 != 41 $\newline$(22,4,5): 49 != 41 $\newline$(22,5,4): 49 != 41 $\newline$(22,6,3): 49 != 41 $\newline$(22,7,2): 49 != 41 $\newline$(22,8,1): 49 != 41 $\newline$

 (23,1,7): 47 != 41 $\newline$(23,2,6): 47 != 41 $\newline$(23,3,5): 47 != 41 $\newline$(23,4,4): 47 != 41 $\newline$(23,5,3): 47 != 41 $\newline$(23,6,2): 47 != 41 $\newline$(23,7,1): 47 != 41 $\newline$

 (24,1,6): 45 != 41 $\newline$(24,2,5): 45 != 41 $\newline$(24,3,4): 45 != 41 $\newline$(24,4,3): 45 != 41 $\newline$(24,5,2): 45 != 41 $\newline$(24,6,1): 45 != 41 $\newline$

 (25,1,5): 43 != 41 $\newline$(25,2,4): 43 != 41 $\newline$(25,3,3): 43 != 41 $\newline$(25,4,2): 43 != 41 $\newline$(25,5,1): 43 != 41 $\newline$

 (26,1,4): 41 = 41 $\newline$(26,2,3): 41 = 41 $\newline$(26,3,2): 41 = 41 $\newline$(26,4,1): 41 = 41 $\newline$

 (27,1,3): 39 != 41 $\newline$(27,2,2): 39 != 41 $\newline$(27,3,1): 39 != 41 $\newline$

 (28,1,2): 37 != 41 $\newline$(28,2,1): 37 != 41 $\newline$

 (29,1,1): 35 != 41 $\newline$In total, there are $4$ ordered triples, so the answer is $\boxed{\textbf{(A)}}$$\newline$$\sim$Imperial Films YT (Unemployed)

 


	\item (\textit{AMC 12 12A 2026}) For what positive integer $n$ is the mean of $1$, $4$, $9$, $16$, $\dots$, $n^{2}$ exactly $32\%$ greater than its median?


 $\textbf{(A)}\   19  \qquad 	\textbf{(B)}\   39  \qquad 	\textbf{(C)}\   49  \qquad 	\textbf{(D)}\   89  \qquad 	\textbf{(E)}\   169$

 \textbf{Answer: }C



\textbf{Solution #1}

Since all answer choices are odd, the median of the sequence is 
 \[\left(\frac{n+1}{2}\right)^2\]

 The mean of the sequence is
 \[\frac{1}{n}\sum_{k=1}^{n} k^2 = \frac{1}{n}\cdot\frac{n(n+1)(2n+1)}{6} = \frac{(n+1)(2n+1)}{6}\]

 Since the mean is $32\%$ greater than the median, we have
 \[\frac{(n+1)(2n+1)}{6} = \frac{132}{100}\left(\frac{n+1}{2}\right)^2\]
 \[200n + 100 = 198n + 198\]
 \[n = \boxed{\textbf{(C) } 49}\]




\textbf{Solution 2 (Divisibility, Answer choices)}

As established in Solution 1, we have the mean and median. Now we write the ratio of of mean to median:
 \[\frac{\text{Mean}}{\text{Median}} = \frac{\frac{(n+1)(2n+1)}{6}}{\frac{(n+1)^2}{4}} = \frac{2(2n+1)}{3(n+1)} = \frac{4n+2}{3n+3}\]

 We are given that this ratio is exactly $32\%$ greater than $1$, which means:
 \[\frac{4n+2}{3n+3}=\frac{132}{100} = \frac{33}{25}\]

 Note that the fraction $\frac{33}{25}$ is in simplest form. For the denominator $3(n+1)$ to reduce to a multiple of $25$ after canceling out any common factors with the numerator, $n+1$ must be a multiple of $25$.

 We test the given answer choices for $n+1$:

 
            \begin{itemize}
        \item (A) $19 + 1 = 20$ (not a multiple of 25)
        \item (B) $39 + 1 = 40$ (not a multiple of 25)
        \item (C) $49 + 1 = 50$ (multiple of 25)
        \item (D) $89 + 1 = 90$ (not a multiple of 25)
        \item (E) $169 + 1 = 170$ (not a multiple of 25)
            \end{itemize}
            Only $n = 49$ satisfies.

 
 \[n = \boxed{\textbf{(C) } 49}\]




\textbf{Solution 3}

Let $n=2k-1.$ The median of the sequence would then be $k^2,$ while the mean would be $\frac{1^2+2^2+3^2+...+(2k-1)^2}{2k-1}=\frac{(2k-1)(2k)(4k-1)}{6(2k-1)}=\frac{(k)(4k-1)}{3}.$ 

 Equating, we get $\frac{(k)(4k-1)}{3}=1.32(k^2)\implies 4k-1=3.96k\implies k=25\implies n=\boxed{\text{(C) }49}$ ~Tacos\_are\_yummy\_1~Minor Edits by Horseneck3

 


	\item (\textit{AMC 12 12A 2026}) What is the value of $(1+i)^{20} + (1-i)^{20}$, where $i = \sqrt{-1}$?


 $\textbf{(A)}\   -2048  \qquad 	\textbf{(B)}\   -1024  \qquad 	\textbf{(C)}\   0  \qquad 	\textbf{(D)}\   1024  \qquad 	\textbf{(E)}\   2048$

 \textbf{Answer: }A



\textbf{Solution 1}

Observe that $(1+i)^2 = 2i$ and $(1-i)^2 = -2i$. $(1+i)^{20} = (2i)^{10} = -1024$ and $(1-i)^{20} = (-2i)^{10} = -1024$. Therefore the answer is $-1024 + -1024 = -2048$, $\fbox{A}.$$\newline$$\sim$ cyrus825

 



\textbf{Solution 2 (Complex Exponentials)}

By Euler's Formula, $a + bi = k(\cos(\theta) + i\sin(\theta)) = ke^{i\theta}$, where $\theta = \tan^{-1}\left(\frac b a\right)$, and $k=\sqrt{a^2+b^2}$ by the Pythagorean Theorem. Therefore, we can express both complex numbers as the following:

 
 \[1  + i = \sqrt{2}e^{i\frac{\pi}{4}}~ \&~  1 - i = \sqrt{2}e^{i\frac{-\pi}{4}}\]

 By DeMoivre's Theorem, $(e^{i\theta})^{n} = e^{in\theta}$. Applying this, we get:

 
 \[(1+i)^{20}  + (1-i)^{20} = 2^{10}(e^{i\frac{\pi}{4}})^{20} + 2^{10}(e^{i\frac{-\pi}{4}})^{20} = 2^{10}\left(e^{i5\pi} + e^{i5\pi}\right) = 2^{11}\left(e^{i\pi}\right) = -2048~\fbox{\textbf{A}}\]$\newline$$\sim$ Fernat123

 



\textbf{Solution 2.1 (Polar Form)}

Alternatively, we can express both numbers directly in polar form without complex exponentials:
 \[1 + i = \sqrt{2}\left(\cos 45^\circ + i\sin 45^\circ\right)\]
 \[1 - i = \sqrt{2}\left(\cos(-45^\circ) + i\sin(-45^\circ)\right)\]

 By De Moivre's Theorem, raising each to the $20^{\text{th}}$ power yields:
 \[(1+i)^{20} = (\sqrt{2})^{20}\left(\cos(20 \cdot 45^\circ) + i\sin(20 \cdot 45^\circ)\right) = 2^{10}\left(\cos 900^\circ + i\sin 900^\circ\right)\]
 \[(1-i)^{20} = (\sqrt{2})^{20}\left(\cos(20 \cdot (-45^\circ)) + i\sin(20 \cdot (-45^\circ))\right) = 2^{10}\left(\cos(-900^\circ) + i\sin(-900^\circ)\right)\]

 Summing the two terms, the imaginary components cancel since $\sin(-900^\circ) = -\sin(900^\circ)$:
 \[(1+i)^{20} + (1-i)^{20} = 2^{10}\left(2\cos 900^\circ\right) = 2^{11}\cos 900^\circ\]

 Because $900^\circ \equiv 180^\circ \pmod{360^\circ}$, we have $\cos 900^\circ = \cos 180^\circ = -1$:
 \[2^{11}(-1) = -2048 \implies \boxed{\textbf{(A) } -2048}\]




\textbf{Solution 3 (Bash)}

First, we compute $(1+i)^5$ by expanding $(1+i)^2$ and $(1+i)^3$:
 \[(1+i)^2=1+2i+i^2=2i\]
 \[(1+i)^3=(1+i)^2(1+i)=2i(1+i)=2i+2i^2=-2+2i\]
 \[(1+i)^5=(1+i)^2(1+i)^3=2i(-2+2i)=-4i+4i^2=-4-4i\]

 Squaring twice gives $(1+i)^{20}$:
 \[(1+i)^{10}=(-4-4i)^2=16+32i+16i^2=32i\]
 \[(1+i)^{20}=(32i)^2=1024i^2=-1024\]

 We do the same for $1-i$:
 \[(1-i)^2=1-2i+i^2=-2i\]
 \[(1-i)^3=(1-i)^2(1-i)=-2i(1-i)=-2i+2i^2=-2-2i\]
 \[(1-i)^5=(1-i)^2(1-i)^3=-2i(-2-2i)=4i+4i^2=-4+4i\]

 Squaring twice:
 \[(1-i)^{10}=(-4+4i)^2=16-32i+16i^2=-32i\]
 \[(1-i)^{20}=(-32i)^2=1024i^2=-1024\]

 Therefore,
 \[(1+i)^{20}+(1-i)^{20}=-1024+(-1024)=\boxed{\textbf{(A) }-2048}.\]




\textbf{Solution 4 (Roots of Unity)}

Let $\omega = \frac{\sqrt{2}}{2} + \frac{\sqrt{2}}{2}i$ be a primitive $8$th root of unity. Then $1 + i = \sqrt{2}\,\omega$ and $1 - i = \sqrt{2}\,\omega^7$, since $\omega^7 = \overline{\omega} = \frac{\sqrt{2}}{2} - \frac{\sqrt{2}}{2}i$.

 Therefore
 \[(1 + i)^{20} + (1 - i)^{20} = (\sqrt{2})^{20}\left(\omega^{20} + \omega^{140}\right) = 2^{10}\left(\omega^{20} + \omega^{140}\right).\]

 Because $\omega^8 = 1$, we reduce the exponents mod $8$. We have $20 \equiv 4 \pmod{8}$ and $140 \equiv 4 \pmod{8}$, so both terms equal $\omega^4$. Since $\omega$ is a primitive $8$th root of unity, $\omega^4$ is the $8$th root of unity directly opposite $1$ on the unit circle, so $\omega^4 = -1$.

 The expression equals $1024(-1 - 1) = \boxed{\textbf{(A)}~-2048}$.

 


 



\textbf{Solution 5 (Conjugates)}

Notice that $1-i=\overline{1+i}$ by the definition of conjugates. Again notice that $\overline{z}^n=\overline{z^n}$ by properties of conjugates, and also that $z+\overline z=2\Re(z)$. Hence we have the following if $z=1+i$:

 
            \begin{itemize}
        \item $(1+i)^{20}+(1-i)^{20}=2\Re((1+i)^{20})$
        \item $\Re((1+i)^{20})=\Re\left(\left(\sqrt 2e^{i\pi /4}\right)^{20}\right)=\Re(2^{10}e^{5i\pi})=-2^{10}$
            \end{itemize}
            Hence our desired value is $2\cdot -2^{10}=-2^{11}=\boxed{-2048\hspace{1mm}\textbf{(A)}}$.

 



\textbf{Solution 6 (Similar to Others)}

Firstly, in order to make the coefficients of 1 and $i$ the same we \textit{must} require that $\cos$ and $\sin$ are the same. This happens when $\theta = \frac{\pi}{4}$. Thus, considering $re^{i \theta}$, we see both expressions evaluate to $\sqrt{2} \cdot e^{i \frac{\pi}{4}}$ and $\sqrt{2} \cdot e^{i \frac{-\pi}{4}}$ respectively. Thus we have

 
 \[(\sqrt{2} e^{i \frac{\pi}{4}})^{20} + (\sqrt{2} e^{i \frac{-\pi}{4}})^{20}\]
 \[2^{10} e^{i5\pi} + 2^{10} e^{-i5\pi}\]
 \[-2^{10} - 2^{10}\]
 \[-1024-1024\]
 \[= \boxed{-2048}.\]



	\item (\textit{AMC 12 12A 2026}) At Euclid High School, 25 students registered to take the AMC 10 and 75 students registered to take the AMC 12. (No student registered to take both.) However, not all 100 students showed up on test day. To compute the percentage of those students who actually took one of the exams, Euclid High School performed the following incorrect two-step process:


 
            \begin{itemize}
        \item Step 1: The school computed both the percentage of AMC 10 registrants who actually took that exam and the percentage of AMC 12 registrants who actually took that exam.
            \end{itemize}
            
            \begin{itemize}
        \item Step 2: The school took the average of the previous two numbers to obtain a final percentage.
            \end{itemize}
            Suppose that Euclid High School's process yielded an answer of $60\%$. What is the minimum number of students who might have actually taken one of the exams?


 $\textbf{(A)}\   30  \qquad 	\textbf{(B)}\   35  \qquad 	\textbf{(C)}\   40  \qquad 	\textbf{(D)}\   50  \qquad 	\textbf{(E)}\   60$

 \textbf{Answer: }C



\textbf{Solution 1:}

Let $x$ be the amount of students who took the AMC 10, and $y$ be the amount of students who took the AMC 12. Clearly, by performing steps one and two, we get $\frac12\left(\frac{x}{25}+\frac{y}{75}\right)=0.6\implies 3x+y=90$. So clearly we need to maximise $x$ if we want to minimise $x+y$. We see that $x\le 25$ by our problem condition, and so $x=25$ yields $y=15$ for a total of $x+y=\boxed{40 \hspace{1mm}\textbf{(C)}}$.

 



\textbf{Solution 1a (skips reasoning of why  needs to be maximised):}

Alternatively, if you really wanna be nitpicky about it, you can set $x+y=k$, then we'd have $2x=90-k\implies x=\frac{90-k}2$. We know that $x\le 25$ so $\frac{90-k}2\le 25\implies 90-k\le 50\implies k\ge 40$ for the answer of $40$.

 



\textbf{Solution 2 (greedy):}

Since the two attendance percentages average to $60\%$, they must sum to $120\%$. To minimize the total number of attendees, we assign $100\%$ to the smaller group and $20\%$ to the larger group, since each percentage point represents fewer students in the smaller group.

 Therefore,

 $1(25)+0.2(75)=\boxed{40 \hspace{1mm}\textbf{(C)}}$




\textbf{Solution 3}

Continue Solution 1 until $3x + y = 90$. Divide by 3. Then, $x + \frac{y}{3} = 30$, and thus $\frac{y}{3} \in \mathbb{Z}$. so set $y = 3y'$ and thus $x + y' = 30$.

 To maximize this, obviously, we want to maximize the number of students. Since $3y' \leq 75 \implies y' \leq 25$ and $x \leq 25$, to minimize we want to minimize $3y'$. This occurs when $x$ is maximized, hence set $x = 25$ and $y' = 5$ in which $y = 15$, hence the answer is $\boxed{40}$.

 


	\item (\textit{AMC 12 12A 2026}) How many ordered pairs $(x,y)$ of real numbers satisfy the following system of equations?
<mjx-container class="MathJax CtxtMenu\_Attached\_0" ctxtmenu\_counter="1" display="true" jax="CHTML" style="font-size: 117.4\%; position: relative;" tabindex="0"><mjx-math aria-hidden="true" class="MJX-TEX" display="true" style="margin-left: 0px; margin-right: 0px;"><mjx-mtable style="min-width: 6.109em;"><mjx-table><mjx-itable><mjx-mtr><mjx-mtd style="text-align: right; padding-right: 0px; padding-bottom: 0.15em;"><mjx-mo class="mjx-n"><mjx-c class="mjx-c7C"></mjx-c></mjx-mo><mjx-msup><mjx-mi class="mjx-i"><mjx-c class="mjx-c1D465 TEX-I"></mjx-c></mjx-mi><mjx-script style="vertical-align: 0.413em;"><mjx-texatom size="s" texclass="ORD"><mjx-mn class="mjx-n"><mjx-c class="mjx-c32"></mjx-c></mjx-mn></mjx-texatom></mjx-script></mjx-msup><mjx-mo class="mjx-n" space="3"><mjx-c class="mjx-c2212"></mjx-c></mjx-mo><mjx-msup space="3"><mjx-mi class="mjx-i"><mjx-c class="mjx-c1D466 TEX-I"></mjx-c></mjx-mi><mjx-script style="vertical-align: 0.413em;"><mjx-texatom size="s" texclass="ORD"><mjx-mn class="mjx-n"><mjx-c class="mjx-c32"></mjx-c></mjx-mn></mjx-texatom></mjx-script></mjx-msup><mjx-mo class="mjx-n"><mjx-c class="mjx-c7C"></mjx-c></mjx-mo><mjx-tstrut></mjx-tstrut></mjx-mtd><mjx-mtd style="text-align: left; padding-left: 0px; padding-bottom: 0.15em;"><mjx-mi class="mjx-n"></mjx-mi><mjx-mo class="mjx-n" space="4"><mjx-c class="mjx-c3D"></mjx-c></mjx-mo><mjx-mn class="mjx-n" space="4"><mjx-c class="mjx-c34"></mjx-c></mjx-mn><mjx-tstrut></mjx-tstrut></mjx-mtd></mjx-mtr><mjx-mtr><mjx-mtd style="text-align: right; padding-right: 0px; padding-top: 0.15em;"><mjx-mfrac><mjx-frac type="d"><mjx-num><mjx-nstrut type="d"></mjx-nstrut><mjx-mrow><mjx-mn class="mjx-n"><mjx-c class="mjx-c31"></mjx-c><mjx-c class="mjx-c36"></mjx-c></mjx-mn><mjx-mi class="mjx-i"><mjx-c class="mjx-c1D465 TEX-I"></mjx-c></mjx-mi></mjx-mrow></mjx-num><mjx-dbox><mjx-dtable><mjx-line type="d"></mjx-line><mjx-row><mjx-den><mjx-dstrut type="d"></mjx-dstrut><mjx-mi class="mjx-i"><mjx-c class="mjx-c1D466 TEX-I"></mjx-c></mjx-mi></mjx-den></mjx-row></mjx-dtable></mjx-dbox></mjx-frac></mjx-mfrac><mjx-tstrut></mjx-tstrut></mjx-mtd><mjx-mtd style="text-align: left; padding-left: 0px; padding-top: 0.15em;"><mjx-mi class="mjx-n"></mjx-mi><mjx-mo class="mjx-n" space="4"><mjx-c class="mjx-c3D"></mjx-c></mjx-mo><mjx-mi class="mjx-i" space="4"><mjx-c class="mjx-c1D465 TEX-I"></mjx-c></mjx-mi><mjx-mi class="mjx-i"><mjx-c class="mjx-c1D466 TEX-I"></mjx-c></mjx-mi><mjx-tstrut></mjx-tstrut></mjx-mtd></mjx-mtr></mjx-itable></mjx-table></mjx-mtable></mjx-math><mjx-assistive-mml display="block" unselectable="on"><math display="block" xmlns="http://www.w3.org/1998/Math/MathML"><mtable columnalign="right left" columnspacing="0em" displaystyle="true" rowspacing="3pt"><mtr><mtd><mo stretchy="false">|</mo><msup><mi>x</mi><mrow data-mjx-texclass="ORD"><mn>2</mn></mrow></msup><mo>-</mo><msup><mi>y</mi><mrow data-mjx-texclass="ORD"><mn>2</mn></mrow></msup><mo stretchy="false">|</mo></mtd><mtd><mi></mi><mo>=</mo><mn>4</mn></mtd></mtr><mtr><mtd><mfrac><mrow><mn>16</mn><mi>x</mi></mrow><mi>y</mi></mfrac></mtd><mtd><mi></mi><mo>=</mo><mi>x</mi><mi>y</mi></mtd></mtr></mtable></math></mjx-assistive-mml></mjx-container>
$\textbf{(A)}\   2  \qquad 	\textbf{(B)}\   4  \qquad 	\textbf{(C)}\   6  \qquad 	\textbf{(D)}\   8  \qquad 	\textbf{(E)}\   10$

 \textbf{Answer: }E



\textbf{Solution 1 (Algebra)}

We get that $\frac{16x}y=xy\implies 16x=xy^2\implies xy^2-16x=0\implies x(y^2-16)=0$. So we have 3 cases to test out.

 
            \begin{itemize}
        \item If $x=0$, then we have $|-y^2|=|y^2|=4$. Since $y^2\ge 0$ by the Trivial Inequality, we must have $y^2=4$ for $y=\pm 2$. This accounts for 2 cases.
        \item If $y^2-16=0$, we have 2 sub-cases. For the positive case of $y=4$, we obtain $|x^2-16|=4\implies x^2-16=\pm 4\implies x^2=16\pm 4=12, 20$ (if the first step is not clear, $|x|=c$ means that $x=c$ or $x=-c$, which is inherently $\pm c$). This provides 4 more cases (a negative and positive $x$ for both numbers). Finally, when $y=-4$ we have $|x^2-16|=4$ which yields 4 more cases.
            \end{itemize}
            Hence we have a total of $2+4+4=\boxed{10\hspace{1mm}\textbf{(E)}}$ ordered pairs.

 



\textbf{Solution 2 (geometric interpretation)}


                    \begin{figure}[H]
                        \centering
                        \includegraphics[width=0.5\textwidth]{Solution/AMC_12_12A_2026_9_Solution_Figure_1.png}
                    \end{figure}
                    

 Note that $\lvert x^2 - y^2 \rvert = 4$ plots two hyperbolae in the plane with vertices $(0, \pm 2)$ and $(\pm 2, 0)$.

 Note that $\frac{16x}{y} = xy$ has solution set $x = 0$ or $y \in \{\pm 4\}$ but not $y = 0$ (not a concern).

 The vertical line $x = 0$ has two points of intersection with the hyperbolae, and the horizontal lines $y \in \{\pm 4\}$ have $4$ points of intersection each with the hyperbolae.

 Hence, the number of real solutions is $2 + 4 + 4 = \boxed{\textbf{(E) }10}$.

 



\textbf{Solution 3 (Logical Cheese)}

Notice that after dividing both sides of the second equation by $x$, we exclude $x=0$ and have two resulting solutions for $y$. Since there are two solutions for $y$ but only one solution for $y^2$, logically, there are two solutions for $x$ under each $y$. Adding back $x=0$ gives us a total of $1+2*2=5$ solutions. Yet, since $x^2 - y^2$ is under an absolute value, we have twice the number of solutions and thus we have $\boxed{\textbf{(E) }10}$.

 


	\item (\textit{AMC 12 12A 2026}) The tree diagram below has nodes labeled with rational numbers. The top node, at level $0$, is $\frac{1}{2}$. Each node $\frac{a}{b}$ sprouts two nodes in the level below it, $\frac{a}{a+b}$ and $\frac{b}{a+b}$. For example, level $1$ has nodes $\frac{1}{1+2} = \frac{1}{3}$ and $\frac{2}{1+2} = \frac{2}{3}$.


 
                    \begin{figure}[H]
                        \centering
                        \includegraphics[width=0.5\textwidth]{Problem/AMC_12_12A_2026_10_Problem_Figure_1.png}
                    \end{figure}
                    


 At what level does the node $\frac{20}{27}$ appear in the diagram?


 

$\textbf{(A)}\   5  \qquad 	\textbf{(B)}\   8  \qquad 	\textbf{(C)}\   12  \qquad 	\textbf{(D)}\   26  \qquad 	\textbf{(E)}\   \text{The node } \dfrac{20}{27} \text{ does not appear in the diagram.}$

 \textbf{Answer: }B



\textbf{Solution 1}

We work backwards. $\frac{20}{27}$ must be below $\frac{7}{20},$ since $\frac{7}{20}$ splits into $\frac{7}{27}$ and $\frac{20}{27}.$ Likewise, we have the chain 
 \[\frac{20}{27}\implies\frac{7}{20}\implies\frac{7}{13}\implies\frac{6}{7}\implies\frac{1}{6}\implies\frac{1}{5},\]where, in the diagram, $\frac{1}{5}$ appears in level $3.$ Thus, $\frac{20}{27}$ appears in level $\boxed{\text{(B) 8}}$ ~Arjun\_Jetti

 Note: You could try to solve your way down, but it would take a lot of time and effort, which you don't want to waste on this competition ~ hellokindsir

 Other Note: However, if you work from left to right, the first node you would work on would be 1/5, so you get lucky there. ~Starlight\_128

 



\textbf{Solution 2}

If $\frac{a}{a+b}$ is emerged from a node, so is $\frac{b}{a+b}$.

 Hence, since $\frac{20}{27}$ comes from a node, so does $\frac{7}{27}$. The result is $\frac{7}{20}$.

 Since $\frac{7}{20}$ is a node, so is $\frac{13}{20}$, so it emerges from $\frac{7}{13}$.

 Since $\frac{7}{13}$ is a node, so is $\frac{6}{13}$, so it emerges from $\frac{6}{7}$.

 Since $\frac{6}{7}$ is a node, so is $\frac{1}{7}$, hence $\frac{1}{6}$.

 From here, it is easy to see we have $\frac{1}{5}, \frac{1}{4}, \frac{1}{3},$ and finally $\frac{1}{2}$.

 Therefore, counting down from 0, it appears on layer $\boxed{8}$.

 



\textbf{Solution 3}

First, let's define neighboring fractions as fractions that have the same number above them (like $1/3$ and $2/3$, or $1/4$ and $3/4$). Each fraction $\frac{a}{b}$ sprouts two fractions $\frac{a}{a+b}$ and $\frac{b}{a+b}$ that add up to $\frac{a+b}{a+b}=1$.

 From this, we see that the neighbor of $\frac{20}{27}$ is $\frac{7}{27}$. Working backwards, the fraction above it is the ratio of the two fractions' numerators, $\frac{7}{20}$.

 We repeat this process:

 
            \begin{itemize}
        \item $7/20$'s neighbor: $13/20$<ul><li>The fraction above: $7/13$</li></ul>
        \item $7/13$'s neighbor: $6/13$<ul><li>The fraction above: $6/7$</li></ul>
        \item $6/7$'s neighbor: $1/7$<ul><li>The fraction above: $1/6$</li></ul>
        \item $1/6$'s neighbor: $5/6$<ul><li>The fraction above: $1/5$</li></ul>
            \end{itemize}
            We see that $1/5$ appears in our diagram in level 3.Therefore, working backwards, $20/27$ appears in level 8, so the answer is $\boxed{\text{(B) 8}}$.

 



\textbf{Solution 3.1}

Each fraction has only one possible fraction it can originate from, so if we start with $\dfrac{20}{27}$, it can only come from $\dfrac{7}{13}$, and that comes from $\dfrac{6}{7}$ which comes from $\dfrac{1}{6}$ which comes from $\dfrac{1}{5}$. (This is just here because I felt the above solution was somewhat confusing)

 



\textbf{Solution 4 (Cheese)}

Notice that answer choice $B$ is exactly 3 more than $A$. The original problem had three pre-established "levels" in the branch, so we can imply that MAA (the test-makers) was trying to trick the people who didn't add the original three "levels". Therefore, the answer is $\boxed{\text{(B) 8}}$.

 



\textbf{Solution 5 (Analyzing Option Choices)}

At level $n$, any node $\frac{a}{b}$ satisfies the denominator bounds:
 \[n + 2 \leq b \leq F_{n+3}\]where $F_3 = 2, F_4 = 3, F_5 = 5, \dots$ are the Fibonacci numbers. Testing $b = 27$ against the options:

 
            \begin{itemize}
        \item \textbf{Eliminate (E):} Since $\gcd(20, 27) = 1$ and $20 < 27$, Euclidean subtraction reduces $\frac{20}{27}$ to $\frac{1}{2}$, so it must appear in the tree.
        \item \textbf{Eliminate (D) $26$:} The minimum denominator at level $26$ is $26 + 2 = 28 > 27$.
        \item \textbf{Eliminate (A) $5$:} The maximum denominator at level $5$ is $F_8 = 21 < 27$.
        \item \textbf{Eliminate (C) $12$:} Reaching $b = 27$ at level $12$ requires staying near the leftmost path ($\frac{1}{14}$), which keeps the numerator far below $20$.
            \end{itemize}
            Thus, $\boxed{8}$ is the only possible answer.

 



\textbf{Solution 6 (No more time left)}

Using the established diagram, we figure the largest possible denominator in level 5 is 21. We immediately eliminate $A$. If we do understand how quickly this pattern moves forward, row 8 should probably cover $\frac{20}{27}$, so we can guess $\boxed{\text{(B) 8}}$. If you are skeptical, you may consider how inflated the numerator and denominator values may be in row 12 and that many AMC 10 questions won't have an answer like $E$ be correct.

 \textbf{Remark} This setup is very similar to the \href{https://en.wikipedia.org/wiki/Calkin\%E2\%80\%93Wilf\_tree}{Calkin-Wilf Tree}, which also contains all rational numbers, but rather than having $\frac{b}{a+b}$, it uses $\frac{a+b}{b}$.

 



\textbf{Solution 7 (Engineer's Induction)}

Looking at the fractions on the far right of the diagram, we can quickly see that the difference in between the numerator and denominator is increasing by $1$. For example, the difference in $\frac{2}{3}$ is $1$, because $3 - 2 = 1$. Likewise, the difference in $\frac{3}{5}$ is 2, and the difference in $\frac{5}{8}$ is 3, because $5 - 3 = 2$ and $8 - 5 = 3$. Let us now use engineer's induction and assume this pattern continues on for this problem. We can also notice that the difference corresponds to the layer/row that the fraction is on by the relationship $n + 1$ (where $n$ is the row). For example, taking the same nodes as before, $\frac{2}{3}$ appears in the second row with a difference of $1$ and therefore $1 + 1 = 2$. Looking at the node $\frac{20}{27}$, we can quickly discern the difference is $7$, and therefore it is on row $7 + 1 = 8$ , making the solution $\boxed{\text{(B) 8}}$



	\item (\textit{AMC 12 12A 2026}) An urn contains $6$ red balls, $4$ black balls, and $4$ white balls. Suppose $7$ balls are drawn at random without replacement and a score is computed as follows:


 - $1$ point is given for each red ball up to a maximum of $3$ points.


 - $1$ point is given for each black ball up to a maximum of $2$ points.


 - $1$ point is given for each white ball up to a maximum of $2$ points.


 For example, drawing $1$ red ball, $4$ black balls, and $2$ white balls would earn a score of $1 + 2 + 2 = 5$. What is the probability of getting the lowest possible score?


 $\textbf{(A)}\   \dfrac{1}{858}  \qquad 	\textbf{(B)}\   \dfrac{5}{3432}  \qquad 	\textbf{(C)}\   \dfrac{1}{312}  \qquad 	\textbf{(D)}\   \dfrac{2}{429}  \qquad 	\textbf{(E)}\   \dfrac{5}{264}$

 \textbf{Answer: }D



\textbf{Solution 1}

The minimum possible score is $4,$ which can occur by drawing $(R,B,W)=\{(6,0,1),(6,1,0),(0,4,3),(0,3,4)\}.$ Considering all balls to be distinct, the first two cases can occur in $2\cdot4\cdot7!$ (ordered) ways - $2$ ways to select black or white, $4$ ways to select the black/white ball, and $7!$ ways to order the $7$ resulting balls. The last two cases total to $2\cdot4\cdot7!.$

 The answer is $\dfrac{2(2\cdot4\cdot7!)}{14\cdot13\cdot12\cdot11\cdot10\cdot9\cdot8}=\boxed{\text{(D) }\dfrac{2}{429}}$ 

 Note: You can think of each case as already having the max amount of balls of a single color and then choosing from the remainder of the balls to get to 4 points. For example: you already have 6 red balls => then you have to choose 1 more ball that gives 1 point. Since all the remaining balls give one point, there are 8 ways(for 8 remaining balls) to do this. If you had 4 black balls, you would have needed 2 more points so you can choose accordingly from the red set of balls.

 -Amon26~Tacos\_are\_yummy\_1~chrislakee: The only hard question on this test imo

 



\textbf{Solution 2 (Casework)}

As explained in the previous solution, the minimum possible score is $4,$ which can occur by drawing $6$ red's with any other color, or $4$ blacks and $3$ whites, or $4$ whites and $3$ black. There are $\binom{14}{7}$ ways to choose 7 balls from a group of 14, and we can now do casework. For the case with 6 reds and 1 of any other color, there are $\binom{6}{6}\cdot \binom{8}{1}=8$, cases. For the other 2 cases, which are symmetric, there are a total of $2\binom{4}{4}\cdot \binom{4}{3}=8$ valid combinations. Thus, we have $\frac{8+8}{\binom{14}{7}}=\frac{16\cdot 7\cdot 6\cdot 5\cdot 4\cdot 3\cdot 2}{14\cdot 13\cdot 12\cdot 11\cdot 10\cdot 9\cdot 8}$. Cancelling terms, we find the fraction simplifies to $\frac{2}{13\cdot 11\cdot 3}=\frac{2}{429}$, so the answer is $\boxed{\textbf{(D)}}$




\textbf{Solution 3}

It is not hard to see that after experimentation that we can achieve a minimal score with:

 
            \begin{itemize}
        \item 6 red, 1 white or black
        \item 3 white and 4 black, or 4 white and 3 black
            \end{itemize}
            Now suppose we line up all 14 balls in a row from left to right. When we draw our first ball in the problem, we pick the leftmost ball, and then we pick the next leftmost ball, and so on. Now let us split into cases:

 
            \begin{itemize}
        \item $\textbf{6 red, 1 white or black}$: We want all 6 red balls and 1 white or black ball to be in the 7 leftmost balls. We have $2 \cdot 7 = 14$ ways to do this. Then for the 7 rightmost balls, we have $\binom{7}{3}=35$ ways of rearranging them and so we have $14 \cdot 35$ ways for this case.
            \end{itemize}
            
            \begin{itemize}
        \item $\textbf{3 white and 4 black, or 4 white and 3 black}$: This is symmetrical to the last case (just swap the 7 leftmost balls with the 7 rightmost balls), so we have $14 \cdot 35$ ways for this case.
            \end{itemize}
            The number of ways to arrange all 14 balls in our row is $\frac{14!}{6! \cdot 4! \cdot 4!}$, and so our desired probability is $\frac{2 \cdot 14 \cdot 35}{\frac{14!}{6! \cdot 4! \cdot 4!}}=\boxed{\textbf{(D)}~\frac{2}{429}}$.

 



\textbf{Solution 4 (BS Cheese)}

As before, the optimal cases are either six red or no red (white and black naturally follow, though this cheese doesn't actually use them). There are likely two distractors to account for the correct value for only one of the cases. If this is true, there is a set of 3 answer choices where two sum to the third, since these cases are mutually exclusive and cover all valid minimums.

 Now, converting each of the fractions to common denominator 3432 to analyze: 
 \[\textbf{(A)}\   \dfrac{4}{3432}  \qquad 	\textbf{(B)}\   \dfrac{5}{3432}  \qquad 	\textbf{(C)}\   \dfrac{11}{3432}  \qquad 	\textbf{(D)}\   \dfrac{16}{3432}  \qquad 	\textbf{(E)}\   \dfrac{65}{3432}\]

 Notably, the only set of 3 choices that have this property are choices $\textbf{(B)}$, $\textbf{(C)}$, and $\textbf{(D)}$. The correct answer corresponds to the sum, $\boxed{\textbf{(D)}\   \dfrac{2}{429}}$.

 Did you actually use this method

 yes, but only to double check lol

 



\textbf{Solution 5 (Very Easy)}

We can see that the minimum amount of points obtained is 4, which can be obtained in two cases. The first case is using all 6 red balls and one of the black or white balls. Since there are 8 black or white balls, there are 8 cases for that. The second case is using only black and white balls. Since we are picking 7 out of the 8 black or white balls, we leave one ball out, and there are 8 ball choices to leave out. Since case 1 has 8 possibilities and case 2 has another 8, and the total possibilities is $\binom{14}{7}$, we get $\boxed{\frac{16}{3432}=\frac{2}{429}}$, which is choice $\boxed{\text{(D)}}$.-Neo

 \textbf{Love this one ~PeetahUndebrethale} 



\textbf{Solution 6 (pure probability, not assumed to be distinct)}

The probability of the first case occurring in one order, 6 red and 1 of either colors, can be found with this:

 $\frac{6}{14} \cdot \frac{5}{13} \cdot \frac{4}{12} \cdot \frac{3}{11} \cdot \frac{2}{10} \cdot \frac{1}{9} \cdot \frac{8}{8}= \frac{1}{3003}$ and then multiply by 7 to get $\frac{1}{429}$, the actual probability of the first case occurring, because there's $\binom{7}{6} = 7$ orders that the events could've happened in.

 The probability of the second case occurring in one order, 4 black and 3 white, can be found with this:$\frac{4}{14} \cdot \frac{3}{13} \cdot \frac{2}{12} \cdot \frac{1}{11} \cdot \frac{4}{10} \cdot \frac{3}{9} \cdot \frac{2}{8} = \frac{1}{30030}$ and then multiply by $\binom{7}{4} = 35$. We also know that the third case, 4 white and 3 black, has the same probability of the second case, so we can do $\frac{1}{30030} \cdot 35 \cdot 2 = \frac{1}{429}$

 Therefore the probability is $\boxed{\textbf{(D)}\   \dfrac{2}{429}}$.

 


	\item (\textit{AMC 12 12A 2026}) Ethan writes down five congruences:

 \[x \equiv  0 \pmod{2}, x \equiv  0 \pmod{3}, x \equiv 1 \pmod{4}, x \equiv 5 \pmod{6}, x \equiv 7 \pmod{12}.\]
How many positive integers less than or equal to $2026$ satisfy at least one of these five congruences?


 $\textbf{(A)}\   1688  \qquad 	\textbf{(B)}\   1856  \qquad 	\textbf{(C)}\   1857  \qquad 	\textbf{(D)}\   2025  \qquad 	\textbf{(E)}\   2026$

 \textbf{Answer: }E



\textbf{Solution 1}

We perform casework on the residues modulo $12$:

 
            \begin{itemize}
        \item The residues $0,2,4,6,8,10$ correspond to $x\equiv 0\pmod{2}$.
        \item The residues $0,3,6,9$ correspond to $x\equiv 0\pmod{3}$.
        \item The residues $1,5,9$ correspond to $x\equiv 1\pmod{4}$.
        \item The residues $5,11$ correspond to $x\equiv 5\pmod{6}$.
        \item The residue $7$ corresponds to $x\equiv 7\pmod{12}$.
            \end{itemize}
            Each residue appears at least once in the list above, hence all integers $x$ satisfy one of these congruences, or $\boxed{\textbf{(E)}}$.

 



\textbf{Solution 2}

This is the first part of AMC 12 so you realize that the problems are not going to be too complex. Write the first few numbers and you can see that they all fit into at least one category. So, you assume that all numbers work, giving you an answer of $\boxed{\text{2026}}$.

 



\textbf{Solution 3}

Notice that all the moduli have a lowest common multiple of $12$, thus if all positive integers between $1$ and $12$ satisfy at least one of the given conditions, then we can simply add $12$ to every single one, and it will satisfy at least the same condition as it does in the base case of integers between $1$ and $12$. For example, the number $5$ satisfies the condition $x \equiv 5 \pmod{6}$, thus if you add $12$, and because $6$ is a factor of $12$, $5+12=17 \equiv 5 \pmod{6}$. Simply testing all numbers from $1$ through $12$ and confirming that every number does in fact satisfy one of these conditions, we can conclude that all positive integers satisfy at least one of these conditions. Therefore, our answer is simply $2026$, or $\boxed{\textbf{E}}$.~Voidling

 


	\item (\textit{AMC 12 12A 2026}) A board has $4$ spaces in a row, which can be filled with numbers. The sides of a fair, $8$-sided die are numbered $1$ through $8$. David rolls the die $4$ times. After each roll, he randomly chooses an empty spot to write the number down on the die. What is the probability that the numbers on the board form a strictly increasing sequence when read from left to right?


 $\textbf{(A)}\   \dfrac{1}{128}  \qquad 	\textbf{(B)}\   \dfrac{5}{432}  \qquad 	\textbf{(C)}\   \dfrac{35}{2048}  \qquad 	\textbf{(D)}\   \dfrac{35}{1024}  \qquad 	\textbf{(E)}\   \dfrac{1}{24}$

 \textbf{Answer: }C



\textbf{Solution}

Notice that this is equivalent to just independently choosing a random number in each of the four places. The number of orderings where the numbers on the board form a strictly increasing sequence is just the number of ways to choose $4$ distinct numbers from the set of $8$ numbers (since each selection corresponds to a unique valid strictly increasing ordering). Therefore, the answer is
 \[\frac{\binom{8}{4}}{8^4}=\frac{8\cdot 7\cdot 6\cdot 5}{8^4\cdot 24}=\frac{70}{2^{12}}=\frac{35}{2048}\]so the answer is $\boxed{\textbf{(C)}}$.

 \textbf{Alternative Solution} The probability that David rolls four different numbers is $\left(\frac{8}{8}\right)\left(\frac{7}{8}\right)\left(\frac{6}{8}\right)\left(\frac{5}{8}\right)$.

 There's only one right order out of 24 orders ($4!$), so the probability that they are in the right order is $\frac{1}{24}$.

 Multiplying, we get $\boxed{\textbf{(C) } \frac{35}{2048}}$.

 \textbf{Alternative Solution 2} We can immediately see that the largest possible first digit is $5$, so we will find the number of correct orderings for each unique first digit $\leq 5$.

 To start, the only valid ordering for the case where the first digit is $5$ is $5, 6, 7, 8$.

 We look at the case where the first number is $4$, and we can see that the remaining $3$ numbers must be selected out of ${5, 6, 7, 8}$. Notice that selecting an ordering whilst satisfying the condition that they must be in strictly increasing order is the same as choosing to exclude $1$ number, as the $3$ numbers that remain only have one order in which they are strictly increasing. Thus the number of orderings where the first digit is $4$ is $\binom{4}{1}$. (The simpler logic is just choosing $3$ of the $4$ numbers and doing $\binom{n}{3}$ for each case, but these two are logically equivalent)

 We can generalize this to the remaining cases of the first digit being $3, 2,$ or $1$. If $3$ is the first digit, then the number of correct orderings is the same as the number of combinations where we exclude $2$ numbers from the $5$ numbers that are bigger than $3$, which is $\binom{5}{2}$.

 Thus, the total number of correct orderings will be $1+\binom{4}{1}+\binom{5}{2}+\binom{6}{3}+\binom{7}{4}=70$.

 The total number of possible orderings is ${8^4}=4096$, so the probability will be $\frac{70}{4096}=\boxed{\textbf{(C) } \frac{35}{2048}}$.

 \textbf{Remark} Answers $\textbf{(B)}$ and $\textbf{(E)}$ are sus because they have denominators that are multiples of 3, but the dice has 8 sides

 


	\item (\textit{AMC 12 12A 2026}) Three identical boxes, each with one open face, are made from $1$-inch-thick wood. The outer dimensions of each box are $a$, $b$, and $c$ inches, where $a < b < c$, and each open face appears on a face with outer dimensions $b \times c$.


 
                    \begin{figure}[H]
                        \centering
                        \includegraphics[width=0.5\textwidth]{Problem/AMC_12_12A_2026_14_Problem_Figure_1.png}
                    \end{figure}
                    


 It happens that if the open faces of the two boxes are held together, then the space inside can exactly contain the (rotated) third box, with no room to spare. What is $a + b + c$?


 $\textbf{(A)}\   20  \qquad 	\textbf{(B)}\   22  \qquad 	\textbf{(C)}\   24  \qquad 	\textbf{(D)}\   26  \qquad 	\textbf{(E)}\   28$

 \textbf{Answer: }C



\textbf{Solution 1}

Joining two of the boxes creates an enclosed space with dimensions $b-2$, $c-2$, and $2(a-1)=2a-2$, where $b-2<c-2<2a-2$. Because the third box with dimensions $a$, $b$, and $c$ must fill this space perfectly, its largest dimension must equal the largest dimension of the enclosed space.
 \[c=2a-2\]

 The same reasoning applies to the smallest dimension, so
 \[a=b-2\]

 The remaining dimension of the box and the enclosed space must match, so
 \[b=c-2\]

 Solving the system of equations,
 \[(a,b,c)=(6,8,10)\implies a+b+c=\boxed{\textbf{(C)}~24}\]




\textbf{Solution 2}

We can directly find that the dimensions of the inner space are $b-2$, $c-2$, and $2(a-1)=2a-2$. Hence, $\{a,b,c\}=\{b-2,c-2,2a-2\}$.

 Notice that $b-2<c-2$, hence the minimal element in the right set is either $b-2$ or $2(a-2)$. In the first case, we have $a=b-2$. Then, either $b=c-2$ or $b=2a-2$. The former yields $a=b-2$, $b=c-2$, and $c=2a-2$, so $(a,b,c)=(6,8,10)$ upon solving. The latter yields $a=b-2$ and $b=2a-2$, but $c=c-2$, impossible. Finally, the second case yields $a=2(a-2)$, so $\{b,c\}=\{b-2,c-2\}$, impossible. Therefore, $(a,b,c)=(6,8,10)$, so the answer is $6+8+10=24$, or $\boxed{\textbf{(C)}}$.

 



\textbf{Solution 3 (Cheese)}

Since the rotated box, when chugged into the $b \times c$ face of the boxes, must have at least 2 dimensions exactly 2 smaller than $b$ and $c$. The 3 side lengths of the box will thus be 3 terms in an arithmetic sequence; the requested sum of the terms is simply thrice the middle term $b$.

 The only answer choice that is a multiple of 3 is $\boxed{\textbf{(C)}~24}$.

 


	\item (\textit{AMC 12 12A 2026}) Suppose that $\sin a = \frac{1}{4}$ and $\sin b = \frac{1}{8}$ with $0 \le b < a \le  \frac{\pi}{2}$. What is the greatest integer less than $\log_{2}\big( \sin(a+b) \cdot \sin(a-b) \big)$?


 

$\textbf{(A)}\   -5  \qquad 	\textbf{(B)}\   -4  \qquad 	\textbf{(C)}\   0  \qquad 	\textbf{(D)}\   1  \qquad 	\textbf{(E)}\   5$

 \textbf{Answer: }A



\textbf{Solution}

We begin by expanding the product inside the logarithm using the sine angle addition and subtraction formulas:
 \[\sin(a + b) = \sin a \cos b + \cos a \sin b\]
 \[\sin(a - b) = \sin a \cos b - \cos a \sin b\]

 Applying the difference of squares identity, the product becomes:
 \[\sin(a + b)\sin(a - b) = \sin^2 a \cos^2 b - \cos^2 a \sin^2 b\]

 Using the fundamental Pythagorean identity $\cos^2 x = 1 - \sin^2 x$, we can rewrite the expression entirely in terms of sine:\begin{align*}\sin^2 a \cos^2 b - \cos^2 a \sin^2 b &= \sin^2 a (1 - \sin^2 b) - (1 - \sin^2 a)\sin^2 b \\&= \sin^2 a - \sin^2 a \sin^2 b - \sin^2 b + \sin^2 a \sin^2 b \\&= \sin^2 a - \sin^2 b\end{align*}


 
Substituting the given values $\sin a = \frac{1}{4}$ and $\sin b = \frac{1}{8}$, we find:
 \[\sin^2 a = \left(\frac{1}{4}\right)^2 = \frac{1}{16}, \quad \sin^2 b = \left(\frac{1}{8}\right)^2 = \frac{1}{64}\]

 Thus, the argument of the logarithm simplifies to:
 \[\sin^2 a - \sin^2 b = \frac{1}{16} - \frac{1}{64} = \frac{4}{64} - \frac{1}{64} = \frac{3}{64}\]

 The expression $E$ is therefore:
 \[E = \log_2 \left( \frac{3}{64} \right)\]

 
To find the maximum integer less than $E$ (the floor of the expression, $\lfloor E \rfloor$), we bound the fraction $\frac{3}{64}$ between two consecutive powers of $2$:
 \[\frac{2}{64} < \frac{3}{64} < \frac{4}{64}\]

 Simplifying these bounding fractions yields:
 \[\frac{1}{32} < \frac{3}{64} < \frac{1}{16}\]

 Since the base-2 logarithm $\log_2(x)$ is a strictly increasing function, we apply it across the inequality:
 \[\log_2 \left( \frac{1}{32} \right) < \log_2 \left( \frac{3}{64} \right) < \log_2 \left( \frac{1}{16} \right)\]

 Recognizing that $\frac{1}{32} = 2^{-5}$ and $\frac{1}{16} = 2^{-4}$, we have:
 \[-5 < E < -4\]

 
The value of the expression lies strictly between $-5$ and $-4$. Therefore, the maximum integer less than the expression is $-5$. 

 This corresponds to choice $\boxed{\text{A}}$.

 Note: Instead of finding the bounds like this, we can also use the property that $\log_2 (\frac{3}{64}) = \log_2 (3) - \log_2(64) = \log_2 (3) - 6$. We see that the floor of $\log_2 (3)$ is $1$, so the answer is $1-6=-5$, or $\boxed{\text{A}}$.

 (And this is emotional damage, cuz at first I got $-5$ for the wrong reason, then I redid it, but this time, I accidentally added them instead of subtracting them despite writing down subtraction and put $-4$)

 



\textbf{Solution 2 (fakesolve)}

The Maclaurin series of $\sin{x} = x - \frac{x^3}{6}...$ and $\arcsin{x} = x + \frac{x^3}{6}...$. For this problem, $\sin{x} = x$ and $\arcsin{x} = x$ suffice (explained at the end).

 Therefore:
 \begin{align*} a = \frac{1}{4}\\ b = \frac{1}{8}\\ \end{align*}


 \begin{align*} \log_2\bigl(\sin(a+b)\cdot\sin(a-b)\bigr) = \log_2(\frac{3}{8} \cdot \frac{1}{8}) = \log_2(\frac{3}{64}) \end{align*}
Note that $-5 < \log_2(\frac{3}{64}) < -4$.Given that the next term in the Maclaurin series expansion of $\sin{x}$ is negative, our estimate for $\sin(a+b)\cdot\sin(a-b)$ is too high. Given that there is only one answer that is less than -4, the answer is thus automatically $\boxed{\textbf{(A) }-5}$.

 Note on why $\arcsin{x} = x$ suffices: We are at most underestimating the value of $a$ by $\frac{1}{384}$ due to the next term, and the value of $b$ by $\frac{1}{3072}$, which would not significantly change the value of the fraction.

 



\textbf{Solution 3 (Engineering)}

We will assume that $\arcsin\frac{1}{4}$ is less than $15$ degrees, and we will use the infamous approximation of $\sin{x}={x}$ for small angles in our solution.

 Replace the expression of $\log_2\bigl(\sin(a+b)\cdot\sin(a-b)\bigr)$ with $\log_2\bigl((\frac{1}{4}+\frac{1}{8})\cdot(\frac{1}{4}-\frac{1}{8})\bigr)$ = $\log_2\bigl(\frac{3}{64}\bigr)$
$\log_2\bigl(\frac{3}{64}\bigr)=\log_2\bigl(3\bigr)-\log_2\bigl(64\bigr)$, which is $($a number between $1$ and $2)-6$.

 Thus our answer choice is $\boxed{\text{A}}$.

 (Note: This is extremely similar to solution 2, except that we don't even consider the Maclaurin series and just make a bunch of assumptions)

 



\textbf{Solution 4}

Note that
 \begin{align*} \cos(2a) &= 1 - 2\sin(a)^2 = 1 - \frac{2}{16} = \frac{7}{8}, \\ \cos(2b) &= 1 - 2\sin(b)^2 = 1 - \frac{2}{64} = \frac{31}{32}. \end{align*}
By product-to-sum,
 \[\sin(a + b)\sin(a - b) = \frac{1}{2}(\cos(2b) - \cos(2a)) = \frac{1}{2}\biggl(\frac{31}{32} - \frac{7}{8}\biggr) = \frac{3}{64}.\]Then,
 \[\log_2\bigl(\sin(a + b) \cdot \sin(a - b)\bigr) = \log_2\biggl(\frac{3}{64}\biggr) = \log_2(3) - \log_2(64) = \log_2(3) - 6,\]which lies in $(-5, -4)$ since $\log_2(3) \in (1, 2)$. Hence, the greatest integer less than $\log_2\bigl(\sin(a + b) \cdot \sin(a - b)\bigr)$ is $\boxed{\textbf{(A) }-5}$.

 



\textbf{Solution 5}

Note that if $\sin a = \frac{1}{4}$, then $\cos a = \frac{\sqrt{15}}{4}$. Similarly, if $\sin b = \frac{1}{8}$, then $\cos b = \frac{\sqrt{63}}{8}$

 The product inside the logarithm can be expanded using the sine angle addition and subtraction formulas:

 
 \[\sin(a + b) = \sin a \cos b + \cos a \sin b\]
 \[\sin(a - b) = \sin a \cos b - \cos a \sin b\]

 Their product can be expressed as a difference of squares: 
 \[\sin^2 a \cos^2 b - \cos^2 a \sin^2 b\]

 Using the given values $\sin a = \frac{1}{4}$ and $\sin b = \frac{1}{8}$, as well as our previously obtained values of $\cos a = \frac{\sqrt{15}}{4}$ and $\cos b = \frac{\sqrt{63}}{8}$, we can evaluate the difference of squares:

 
 \[\left(\frac{1}{4}\right)^2 \cdot \left(\frac{\sqrt{63}}{8}\right)^2 - \left(\frac{\sqrt{15}}{4}\right)^2 \cdot \left(\frac{1}{8}\right)^2 = \frac{63}{1024} - \frac{15}{1024} = \frac{48}{1024} = \frac{3}{64}\]

 Our desired value is now the greatest integer less than or equal to $\log_2 \left( \frac{3}{64} \right)$. Using the quotient rule of logarithms, this expression can be simplified to $\log_2 3 - \log_2 64$, and then to $\log_2 3 - 6$ since $2^6 = 64$. 

 Since $2 < 3 < 4$ and $\log_2$ is increasing, $\log_2 3 \in (1, 2)$. It follows that $(\log_2 3 - 6) \in (-5, -4)$, and the greatest integer less than or equal to this is $\boxed{\textbf{(A) }-5}$.

 



\textbf{Solution 6}

Recall the identity

 
 \[\sin(a + b)\sin(a - b) = \sin^2 a - \sin^2 b.\]

 Thus, we must evaluate $\left\lfloor \log_2\left(\left(\frac{1}{4}\right)^2 - \left(\frac{1}{8}\right)^2\right) \right\rfloor = \left\lfloor \log_2 \frac{3}{64} \right\rfloor$. Since $\frac{1}{32} < \frac{3}{64} < \frac{1}{16}$, we have $-5 < \log_2 \frac{3}{64} < -4$, so the answer is $\boxed{\textbf{(A) }-5}$.

 


	\item (\textit{AMC 12 12A 2026}) Gerry is playing the Triples Game. Her task is to partition the $3 \times 3$ grid of squares shown on the left into three sections, each containing three adjoining squares. (Three squares are adjoining if each shares an edge with at least one of the other two.) One point will be awarded for each section containing more $\textsf{X}$'s than $\textsf{O}$'s. Shown on the right is a sample partition that will earn one point in total. A partition that earns two or more points will win the game. What fraction of the possible partitions of this particular grid will result in a win for Gerry?


 
                    \begin{figure}[H]
                        \centering
                        \includegraphics[width=0.5\textwidth]{Problem/AMC_12_12A_2026_16_Problem_Figure_1.png}
                    \end{figure}
                    
$\textbf{(A)}\   \dfrac{3}{10}  \qquad 	\textbf{(B)}\   \dfrac{1}{3}  \qquad 	\textbf{(C)}\   \dfrac{2}{5}  \qquad 	\textbf{(D)}\   \dfrac{4}{9}  \qquad 	\textbf{(E)}\   \dfrac{1}{2}$

 \textbf{Answer: }C



\textbf{Solution 1}

There are four corner squares but only three sections, so some section contains two corners. A connected section of three squares containing two corners must be a complete outer row or column.

 Removing it leaves a $2\times3$ rectangle, which splits into either two straight sections or two L-shaped sections. There are $2$ all-straight partitions (rows or columns). Otherwise, choose the single straight section on one of the $4$ sides, then split the remaining rectangle into two L-shapes in $2$ ways. Thus there are
 \[2+4\cdot2=10\]partitions in total.

 Since there are only $4$ X's, a win requires exactly $2$ X's in each of two sections. All $10$ partitions are shown below; exactly $4$ win.

 
                    \begin{figure}[H]
                        \centering
                        \includegraphics[width=0.5\textwidth]{Solution/AMC_12_12A_2026_16_Solution_Figure_1.png}
                    \end{figure}
                    

 Therefore, the requested fraction is
 \[\frac{4}{10}=\boxed{\frac25},\]We choose $\boxed{\textbf{(C)}}$.

 



\textbf{Solution 2}

Consider the section that includes the bottom-right square. Note that it will always have two O's, so whatever section contains it will be invalid. We casework over all the possible sections, and look at how we can split the remaining 6 tiles for each. This makes it easier to count, as now both the other sections must have more X's than O's to win, and we are sure that we have covered all the cases. Proceed as Solution 1.

 



\textbf{Solution 3 (similar to Solution 2)}

Consider the section that includes the bottom-right square. Note that it will always have two O's.  If you really want to know how you can immediately get 10, you can split it into casework that either there are 3 rows which satisfy this, one row which satisfies this (and the rest are weird L's), three columns which satisfy this, and one column which satisfies this. To create three rows and three columns, it's obviously just 1 for both of them, which is 2 in total. For the L's case the one row can either be the first row or the third row by logic and the L's can have 2 different positions to be in to satisfy the 2x3 remaining rectangle. Therefore it is 2x2, which is 4, and by columns, the answer remains the same. Therefore, the number of permutations in total is 10. It's better to see the answer choices to make sure we're correct. the 9 in 4/9 is awfully close to our 10, but under further inspection none of our cases intersect. That too, we got 4 for both of them, meaning we eventually get $\boxed{\textbf{(C)}}$. Proceed as above.

 



\textbf{Solution 4}

Consider the section that includes the bottom-right square. Note that it will always have two O's, so whatever section contains it will be invalid. We casework over all the possible sections, and look at how we can split the remaining 6 tiles for each. This makes it easier to count, as now both the other sections must have more X's than O's to win, and we are sure that we have covered all the cases. Proceed as Solution 1.

 GoldemGem4848: solution 2 achieved mitosis

 



\textbf{Solution 5}

First, we count the number of total partitions. Consider the "$l$" shape formed by three squares. We see that if we fix any "$l$" that contains the center square, the rest is uniquely determined. There are 8 such "$l$"s, but we have to also account for the case where there are 3 vertical lines and 3 horizontal lines. These account for 10 partitions in total.

 Now, we count the number of winning partitions. Let the square directly to the left of the center square be called the penguin square.

 Case 1: Center and penguin square are together in a winning partition. We can see it either has to be horizontal or containing the bottom-left square, and the rest are uniquely determined, for a total of 2 cases.

 Case 2: Center and penguin square are disjoint in a partition. We see that the center square has to be in the partition with the top-right square, and then the rest are uniquely determined, for a total of 2 cases.

 $\frac{2+2}{10} = \boxed{\textbf{(C)}~\frac{2}{5}}$.

 



\textbf{Solution 6}

Consider the section that includes the bottom-right square. Note that it will always have two O's, so whatever section contains it will be invalid. We casework over all the possible sections, and look at how we can split the remaining 6 tiles for each. This makes it easier to count, as now both the other sections must have more X's than O's to win, and we are sure that we have covered all the cases. Proceed as Solution 1.~Moonwatcher22

 



\textbf{Solution 7}

For each section that gains a point, it must have atleast 2 X's as if it only has maximum one X, it will have at least 2 O's and not count. So that means (a) max 2 points can be gained and (b) there must be one section with only O's as otherwise there are only maximum 3 X's and only one point can be gained. There are 2 ways this can be achieved: either you use the bottom row of O's and then you have 3 ways, either 3 rows, or the bottom row and 2 L-shapes which can be oriented in 2 ways. The 3 rows and one orientation of the L-shapes work, so we get 2 cases. The other case is the bottom L-shape of 3 O's and any of the 2 ways from here works. So in total there are 4 ways that work, and to count the total ways, we have 2 ways of only 3 rows or 3 columns, and then we have the ones with the L-shapes: one line in any direction gives 2 orientations of the L-shapes and there are 2 rows and 2 columns that allow L-shapes, which makes $2\cdot4=8$ ways so in total we have $2+8=10$ ways so our answer is $\frac{4}{10}$ which simplifies to $\boxed{\text{(C) }\frac{2}{5}}$.~keyo

 


 


	\item (\textit{AMC 12 12A 2026}) An engineer is designing a closely packed field of $7$ wind turbines for electricity production. Each turbine, whose center is marked by a labeled dot in the figure, occupies one of $7$ congruent elliptical regions. Turbines in the same row have centers that are $4$ units apart, and turbines in the same column have centers that are $9$ units apart, as shown in the figure below.


 
                    \begin{figure}[H]
                        \centering
                        \includegraphics[width=0.5\textwidth]{Problem/AMC_12_12A_2026_17_Problem_Figure_1.png}
                    \end{figure}
                    


 The major axis of each ellipse is $m\sqrt{p}$, where $m$ and $p$ are positive integers and $p$ is not divisible by the square of any prime. What is $m + p$?


 $\textbf{(A)}\   6  \qquad 	\textbf{(B)}\   8  \qquad 	\textbf{(C)}\   9  \qquad 	\textbf{(D)}\   10  \qquad 	\textbf{(E)}\   12$

 \textbf{Answer: }A



\textbf{Solution 1}

Let $L$ be the major axis length. Each ellipse has horizontal semiaxis $2$, and adjacent rows are $\frac{9}{2}$ apart. Place the middle center at $O=(0,0)$; then the upper-right center is $A=(2,\frac{9}{2})$.

 Let $T$ be their point of tangency. A half-turn about the midpoint of $OA$ swaps the two ellipses, so it fixes their unique common point $T$. Thus $T$ is the midpoint of $OA$.

 Dilating the middle ellipse by a factor of $2$ about $O$ sends $T$ to $A$. The resulting ellipse, drawn with a blue dashed curve, has semiaxes $4$ and $L$.

 
                    \begin{figure}[H]
                        \centering
                        \includegraphics[width=0.5\textwidth]{Solution/AMC_12_12A_2026_17_Solution_Figure_1.png}
                    \end{figure}
                    

 Since $A=(2,\frac{9}{2})$ lies on this ellipse,
 \[\frac{2^2}{4^2}+\frac{(9/2)^2}{L^2}=1 \quad\Longrightarrow\quad L^2=27 \quad\Longrightarrow\quad L=3\sqrt3.\]Therefore $m+p=3+3=\boxed{\text{(A)}\ 6}$.

 



\textbf{Solution 2}


                    \begin{figure}[H]
                        \centering
                        \includegraphics[width=0.5\textwidth]{Solution/AMC_12_12A_2026_17_Solution_Figure_2.png}
                    \end{figure}
                    Let $O$ be the origin, at $(0,0)$. Also, let the major axis of an ellipse, which is what we're looking for, be $a$. Then the equation of the ellipse centered at $O$ is
 \[\frac{x^2}{4} + \frac{y^2}{(\frac{a}{2})^2}=1.\]

 Notice that if we could find one pair that lies on the ellipse, we could plug it into the equation and solve for $a$.

 Let's focus on the ellipse centered at $A$. It has coordinates $(2, \frac{9}{2})$. Notice that the midpoint of $O$ and $A$ ($M$ in the above diagram) lies on the central ellipse. The coordinates of that midpoint are $(1, \frac94)$ (by the midpoint formula). We now have a point that lies on the ellipse. Plugging it in, we get

 
 \[\frac{1^2}{4} + \frac{(\frac94)^2}{(\frac{a}{2})^2}=1.\]
 \[a^2 + 81 = 4a^2\]
 \[a^2 = 27.\]

 Thus, $a = 3\sqrt3$, and our answer is $\boxed{\text{(A) }6}.$

 Hi, would you want a diagram added to your solution?

 Hi, yes.

 Hi, I have added the diagram (though i'm not the person who asked the original question) ~sujoybec51

 


 



\textbf{Solution 2.1 (Implicit Differentiation)}

Only do this if you forgot midpoint formula, or if you love complicating things while doing Solution 2.

 The equation of the ellipse with center $(0,0)$ is 
 \[\frac{x^2}{4} + \frac{y^2}{b^2} = 1.\]

 The equation for the ellipse centered at $(2, \frac{9}{2})$ is
 \[\frac{(x - 2)^2}{4} + \frac{\left(y - \frac{9}{2}\right)^2}{b^2} = 1.\]

 Let $T(x, y)$ be the point of tangency between ellipses $O$ and $A$. Because they share the same tangent slope $\frac{dy}{dx}$ at $T$, differentiating both equations implicitly with respect to $x$ gives
 \[\frac{dy}{dx} = -\frac{b^2 x}{4y} = -\frac{b^2(x - 2)}{4\left(y - \frac{9}{2}\right)}.\]

 Simplifying:
 \[\frac{x}{y} = \frac{x - 2}{y - \frac{9}{2}}\]
 \[y = \frac{9}{4}x.\]

 To find where this line intersects the ellipses, set the two ellipse equations equal at $(x, y)$ and simplify:
 \[\frac{x^2}{4} + \frac{y^2}{b^2} = \frac{(x - 2)^2}{4} + \frac{\left(y - \frac{9}{2}\right)^2}{b^2}\]
 \[(1 - x) + \frac{\frac{81}{4} - 9y}{b^2} = 0.\]

 Substituting $y = \frac{9}{4}x$ gives
 \[(1 - x)\left(1 + \frac{81}{4b^2}\right) = 0.\]

 Since $1 + \frac{81}{4b^2} > 0$, by Zero Product Property, we know that $1 - x = 0$, from which we obtain $x = 1$, and therefore
 \[y = \frac{9}{4}.\]

 The rest is the same as Solution 2. You should get$m = 3$ and $p = 3$, so
 \[m + p = \boxed{\textbf{(A) }6}.\]




\textbf{Solution 3 (Easiest)}

We can view the given figure of ellipses as a vertical stretch of this figure of seven tangent circles of radius $2$ by some scale factor: 

 
                    \begin{figure}[H]
                        \centering
                        \includegraphics[width=0.5\textwidth]{Solution/AMC_12_12A_2026_17_Solution_Figure_3.png}
                    \end{figure}
                    

 By connecting three centers as shown, we create an isosceles triangle with equal legs of length $4$ and deduce $h=4\sqrt3.$ Since this vertical distance was given as $9$ in the problem, we conclude that the stretch factor between this diagram and the given one is $\frac{9}{4\sqrt3}.$ 

 Thus, the major axis of one of the ellipses is simply the diameter of one of these circles multiplied by the stretch factor, or $m\sqrt{p}=(4)\left(\frac{9}{4\sqrt3}\right)=3\sqrt{3}\implies m+n=\boxed{\text{(A) } 6}$




\textbf{Solution 4 (Ruler Cheese)}

You can ruler bash. Very risky but works. Edit: By ruler, we get the length of one major axis to be around $5.125$, which is very close to $3 \sqrt3$, so $3+3=$ $\boxed{\text{(A) }6}.$

 It would be so orz if someone actually got the answer by this :sob:

 \textbf{Note} Yay, a fellow cheese enthusiast!

 A way to justify the radical approximation 5.125 is as follows:

 Due to the hexagonal nature (and rightfully so, to maximize electricity, "hexagons are the bestagons") of the turbines, the simplified number under the radical is almost certainly $p=3$. From there, you can 3 from the answer choices to obtain $m$: 
 \[\textbf{(A)}~m=3 \qquad \textbf{(B)}~m=4 \qquad \textbf{(C)}~m=6 \qquad \textbf{(D)}~m=7 \qquad \textbf{(E)}~m=9\]

 Also, here are some calculations to show the diagram is intended to be to scale for the ruler cheese:

 You can verify that the ratio of the small to the long line segment on your test sheet is indeed approximately 4:9. For example, the measurement on my test sheet, to the rounding error in my 0.1 cm ruler, is 1.4 cm:3.4 cm; using the same ruler, the measurement on the diagram on this webpage above at 100\% zoom is 2.7 cm: 6.7 cm.

 Since the scale on both axes is similar to the given values with negligible error, we conclude the graph is almost certainly to scale.

 



\textbf{Solution 5 (Transformations)}

We apply the transformation to the diagram 
 \[T(x, y) \implies \left( \frac{x}{a}, \frac{y}{b} \right)\]So that the center turbine maps to a unit circle.

 
                    \begin{figure}[H]
                        \centering
                        \includegraphics[width=0.5\textwidth]{Solution/AMC_12_12A_2026_17_Solution_Figure_4.png}
                    \end{figure}
                    Since each of the outer circles are contacting each other and the center circle, hexagon ABCDEF is a regular hexagon with a side length of 2. On the original hexagon, the distance $\overline{AE}$ is 9. Under $T(x, y)$, the distance AE becomes $\frac{9}{a}$.The distance between AE is also equal to 

 $\overline{AE}= 2*2\sin(60) = 2\sqrt{3}$. 

 Therefore, 

 $\frac{9}{a} = 2\sqrt{3} \implies a= \frac{3\sqrt{3}}{2}$ 

 Hence the major axis is $2a = 3\sqrt{3}$. The answer is $3 + 3 = \boxed{\text{(A) }6}.$



	\item (\textit{AMC 12 12A 2026}) A sphere with center $O$ and radius $9$ cuts three mutually perpendicular planes that meet at point $P$, forming three circular cross-sections on the planes, which have radii $7$, $4$, and $3$. What is $OP$?


 $\textbf{(A)}\   10  \qquad 	\textbf{(B)}\   11  \qquad 	\textbf{(C)}\   12  \qquad 	\textbf{(D)}\   13  \qquad 	\textbf{(E)}\   14$

 \textbf{Answer: }D



\textbf{Solution 1}

Every point on the three cross-sections also lies on the sphere. Therefore, each of those points is $9$ away from $O$ from the definition of the sphere. Let the centers of the circles with radii $7$, $4$, and $3$ be $A$, $B$, and $C$, respectively. By the Pythagorean Theorem, 
 \[OA = \sqrt{9^2 - 7^2} = \sqrt{32}.\]
 \[OB = \sqrt{9^2 - 4^2} = \sqrt{65}.\]
 \[OC = \sqrt{9^2 - 3^2} = \sqrt{72}.\]

 Assume that the planes are the $xy$ plane, the $xz$ plane, and the $yz$ plane, all of which are mutually perpendicular. The distance from an arbitrary point to each plane is equal to the $z$, $y$, and $z$ coordinates of that point. Observe that $OA$, $OB$, and $OC$ are the distances from $O$ to each of the planes. Then the distance between $O$ and the origin, or point $P$, is equal to $\sqrt{x^2+y^2+z^2}$, where $x$, $y$, and $z$ are the coordinates of $O$. By our reasoning above, the coordinates of $O$ are $(\sqrt{32},\sqrt{65},\sqrt{72})$. Plugging in, we find that 
 \[OP = \sqrt{(\sqrt{32})^2+(\sqrt{65})^2+(\sqrt{72})^2} = \sqrt{32+65+72} = \sqrt{169} = \boxed{13},\]so the answer is $\boxed{(D)\, 13}$.~gb1falcon

 


	\item (\textit{AMC 12 12A 2026}) How many values of $x$ between $0$ and $2\pi$, inclusive, satisfy the equation

 \[\sin x + \sin(2x) + \sin(3x) = 1 + \cos x + \cos(2x)\ ?\]
$\textbf{(A)}\   1  \qquad 	\textbf{(B)}\   2  \qquad 	\textbf{(C)}\   4  \qquad 	\textbf{(D)}\   6  \qquad 	\textbf{(E)}\   8$

 \textbf{Answer: }D



\textbf{Solution 1}

Use the identity $\sin(3x)=3\sin x-4\sin^3 x$ as well as the double-angle identities to get
 \[4\sin x-4\sin^3x +2\sin x\cos x=1+\cos x+2\cos^2x-1=\cos x+2\cos^2 x\]We know that $4\sin x-4\sin^3x=4\sin x(1-\sin^2x)=4\sin x \cos^2 x$. Therefore,
 \[4\sin x\cos^2 x+2\sin x\cos x=\cos x+2\cos^2 x\]We may then extract $\cos x=0$. Otherwise, dividing by $\cos x$ yields
 \[4\sin x\cos x+2\sin x=1+2\cos x\implies (2\cos x+1)(2\sin x-1)=0\]which implies that $\cos x=-\frac{1}{2}$ or $\sin x=\frac{1}{2}$. All three sets of solutions yield exactly two solutions each, all unique, hence there are $6$ total solutions, or $\boxed{\textbf{(D)}}$.

 



\textbf{Solution (Complex Numbers)}

Let $z=e^{ix}$ and $S=1+z+z^2$. The given equation is
 \[\operatorname{Im}(zS)=\operatorname{Re}(S).\]Since $|z|=1$, we have $\overline S=S/z^2$ and $\overline{zS}=S/z^3$. Hence
 \[\frac{zS-S/z^3}{2i}=\frac{S+S/z^2}{2}.\]Multiplying by $2iz^3$ and factoring gives
 \[0=S\bigl(z^4-1-iz(z^2+1)\bigr) =(z^2+z+1)(z^2+1)(z^2-iz-1).\]The three factors give
 \[z=\frac{-1\pm i\sqrt3}{2}, \qquad z=\pm i, \qquad z=\frac{i\pm\sqrt3}{2}.\]These are six distinct points on the unit circle, none equal to $1$. Each therefore gives exactly one $x$ in $[0,2\pi]$, so the answer is $\boxed{6}$, we choose $\boxed{\textbf{(D)}}$.

 


 



\textbf{Solution 2}

We can use the sum-to-product identity 
 \[\sin(a)+\sin(b)=2\sin\left(\frac{a+b}2\right)\cos\left(\frac{a-b}2\right)\] on $\sin(x)+\sin(3x)$ to obtain $2\sin(2x)\cos(x)$ . Hence we have 
 \[2\sin(2x)\cos(x)+\sin(2x)=1+\cos(x)+\cos(2x)\implies \\  \sin(2x)(2\cos(x)+1)=2\cos^2(x)+\cos(x)=\cos(x)(2\cos(x)+1).\] At this point, we can see that the solutions to $2\cos(x)+1=0$ satisfy this equation, which will yield 2 solutions on $[0, 2\pi]$. Now we cancel and get $\sin(2x)=\cos(x)\implies 2\sin(x)\cos(x)=\cos(x)$. So the solutions to $\cos(x)=0$ also satisfy our original equation, which yields another 2 solutions. Finally, $2\sin(x)=1$ yields another 2 solutions, for a total of $2+2+2=\boxed{6\hspace{1mm}\textbf{(D)}}$ solutions.

 



\textbf{Solution 3 (Intuition)}

Write $\sin x = \sin(2x - x)$ and $\sin 3x = \sin(2x + x)$. Expanding gives
 \[\sin(2x - x) = \sin 2x\cos x - \cos 2x\sin x,\]
 \[\sin(2x + x) = \sin 2x\cos x + \cos 2x\sin x.\]Adding this has the $\cos 2x\sin x$ terms cancel, so
 \[\sin x + \sin 3x = 2\sin 2x\cos x.\]The LHS therefore becomes
 \[\sin 2x + 2\sin 2x\cos x = \sin 2x\,(1 + 2\cos x).\]

 For the RHS, use $\cos 2x = 2\cos^2 x - 1$ to get
 \[1 + \cos x + 2\cos^2 x - 1 = \cos x\,(1 + 2\cos x).\]

 Setting the two sides equal and factoring out $1 + 2\cos x$ gives
 \[(1 + 2\cos x)(\sin 2x - \cos x) = 0.\]Since $\sin 2x - \cos x = 2\sin x\cos x - \cos x = \cos x\,(2\sin x - 1)$, the equation factors out into
 \[(1 + 2\cos x)\cos x\,(2\sin x - 1) = 0.\]

 Now all that’s left is to count solutions on $[0, 2\pi]$:

 
            \begin{itemize}
        \item $\cos x = -\frac{1}{2}$ gives $x = \frac{2\pi}{3}, \frac{4\pi}{3}$.
        \item $\cos x = 0$ gives $x = \frac{\pi}{2}, \frac{3\pi}{2}$.
        \item $\sin x = \frac{1}{2}$ gives $x = \frac{\pi}{6}, \frac{5\pi}{6}$.
            \end{itemize}
            Thus the answer is $\boxed{\textbf{(D)}~6}$.

 



\textbf{Solution 4 (Graphing)}


                    \begin{figure}[H]
                        \centering
                        \includegraphics[width=0.5\textwidth]{Solution/AMC_12_12A_2026_19_Solution_Figure_1.PNG}
                    \end{figure}
                     Instead of using identities or other tricks, we can simply graph the equations $y=\sin x+\sin\left(2x\right)+\sin\left(3x\right)$ and $y=1+\cos\left(x\right)+\cos\left(2x\right)$. After counting the intersections  we get the answer $\boxed{\textbf{(D)}~6}$.

 Note: Although we don't have Desmos on the test, graphing it using every theta value we know trig for ex: ($\frac{\pi}{6}$, $\frac{\pi}{3}$, is good enough for a reasonably accurate graph. ~AncientKarnivore

 


	\item (\textit{AMC 12 12A 2026}) How many polynomials $f(x)$ of degree $3$ satisfying

 \[f\left( 3+\sqrt{5} \right) = f\left( 3 - \sqrt{5} \right)\]
are there that have leading coefficient equal to $1$ and all of whose coefficients are integers between $-10$ and $10$, inclusive?


 $\textbf{(A)}\   21  \qquad 	\textbf{(B)}\   42  \qquad 	\textbf{(C)}\   63  \qquad 	\textbf{(D)}\   84  \qquad 	\textbf{(E)}\   105$

 \textbf{Answer: }D



\textbf{Solution 1}

Let the polynomial be $f(x)=x^3+ax^2+bx+c$ as the leading coefficient is 1. Substituting yields
 \[\left(3-\sqrt{5}\right)^3+a\left(3-\sqrt{5}\right)^2+b\left(3-\sqrt{5}\right)+c= \left(3+\sqrt{5}\right)^3+a\left(3+\sqrt{5}\right)^2+b\left(3+\sqrt{5}\right)+c\]Expanding yields
 \[(27-27\sqrt{5}+45-5\sqrt{5})+a(9-6\sqrt{5}+5)+b(3-\sqrt{5})+c\]
 \[=(27+27\sqrt{5}+45+5\sqrt{5})+a(9+6\sqrt{5}+5)+b(3+\sqrt{5})+c\]Mass cancelling yields
 \[32\sqrt{5}+6\sqrt{5}a+\sqrt{5}b=0\implies 6a+b=-32\]Notice that since $-10\le b\le 10$, we get
 \[-42= -32-10\le -32-b\le -32+10=-22\]Hence,
 \[-7\le a\le -\frac{11}{3}\implies -7\le a\le -4\]Substituting yields the valid solutions $(a,b)=(-4,-8),(-5,-2),(-6,4),(-7,10)$. Noting that $c$ cancels and can thus take on any of the $21$ integers from $-10$ to $10$, inclusive, our answer is $21\cdot 4=84$, or $\boxed{\textbf{(D)}}$.

 



\textbf{Solution 2}

Let $f\left(3+\sqrt 5\right)=f\left(3-\sqrt 5\right)=c$. Now define a new function $g(x)\coloneqq f(x)-c$. Clearly $g$ is still a cubic, and the 2 of the roots of $g$ will be $r_1, r_2=3+\sqrt 5, 3-\sqrt 5$. Hence $g(x)=(x-d)(x-r_1)(x-r_2)+c$. We work backwards to get 
 \[x=3\pm\sqrt 5\implies x-3=\pm\sqrt 5\implies (x-3)^2=5\implies x^2-6x+9=5\implies x^2-6x+4=0.\] So $g(x)=(x-d)(x^2-6x+4)+c$. Now we just need to apply our given condition. Expanding we have that $g(x)=x^3-(6+d)x^2+(4+6d)x-4d+c$. The most restrictive inequality would be applied to the coefficient of $x$, to yield 
 \[-10\le 4+6d\le 10\implies -14\le 6d\le 6\implies -\frac73\le d\le 1.\]So $d\in\{-2, -1, 0, 1\}$. We can choose $c$ freely as long as $|-4d+c|\le 10$, which gives 21 valid $c$ for each case. Hence there are $4\cdot 21=\boxed{84\hspace{1mm}\textbf{(D)}}$ such functions. 

 Also note that $c$ can have 21 choices because it cancels out in the initial equality of $f\left(3+\sqrt 5\right)=f\left(3-\sqrt 5\right)$~Amon26

 Note: If you are stuck on which answer to choose, once you have gotten to the restrictions on $d$, you can clearly note that the answer has to be divisible by 4, in which $\textbf{(D)}$ is the only option.

 



\textbf{Solution 3}

First let $x_1=3+\sqrt5$ and $x_2=3-\sqrt5.$ 

 Next, let $f(x)=x^3+ax^2+bx+c$ since $f$ is monic. From the problem statement, 
 \[x_1^3+ax^2_1+bx_1+c=x_2^3+ax_2^2+bx_2+c.\]

 Clearly, the value of $c$ is irrelevent since subtracting $c$ from both sides retains the equation. So, there are $21$ choices. (A check is that the answer choices are multiples of $21$)

 Now, substitute for $x_1$ and $x_2.$
$x_1^2=14+6\sqrt5$ and $x_2^2=14-6\sqrt5$ next$x_1^3=72+32\sqrt5$ and $x_2^3=72-32\sqrt5$.

 So, substituting one gets 
 \[(72+32\sqrt5)+a(14+6\sqrt5)+b(3+\sqrt5)=(72-32\sqrt5)+a(14-6\sqrt5)+b(3-\sqrt5).\]

 Cancelling terms and rearranging gives 

 
 \begin{align*} 64\sqrt5+12a\sqrt5+2b\sqrt5&=0 \Rightarrow \\ 32+6a+b&=0. \end{align*}


 Take the equation$\pmod 6$ yields $2+b\equiv 0 \pmod 6$ so $b \equiv 4 \pmod 6$ and $b=-8, -2, 4, 10.$

 Substituting yields $a=-4, -5, -6, -7$ respectively. 

 Thus, since there are only four pairs $(a, b)$ and $21$ choices of $c$ for each pair there are $4\cdot21=\boxed{84\hspace{1mm}\textbf{(D)}}$ total pairs.

 Note this solution follows similar to solution $1$ but finishes with a quick modulo inspection.Some motivation of using the form $x^3+ax^2+bx+c$ rather than a function like $(x-r_1)(x-r_2)(x-r_3)$ is because the problem references the coefficients.

 



\textbf{Solution 4 (Similar to Solution 2)}

Let $f\left( 3+\sqrt{5}\right) = f\left( 3-\sqrt{5}\right) = c$ for some constant $c$. Consider the function $f(x)-c$. This function has roots at $3 \pm \sqrt{5}$ and is monic (leading coefficient $1$), so it is divisible by the monic quadratic which also has these two roots.

 
We rewrite $3 \pm \sqrt{5}$ as $\frac{6\pm \sqrt{20}}{2}$. By the quadratic formula, the coefficient of the linear term is $-6$ and the leading coefficient is $1$. We can then find that the constant is $4$. We can now set $f(x)-c=(x^2-6x+4)(x-k)$.

 
We expand the right hand side and rearrange to get $f(x)=x^3-(k+6)x^2+(6k+4)x-4k+c$. We bound the linear coefficient $6k+4$ between $[-10, 10]$ to get $k=-2, -1, 0, 1$. For any of these four values of $k$, there are $21$ values of $c$ such that the constant is in $[-10, 10]$. Furthermore, the quadratic coefficient is always in this range for all choices of $k$. 

 
To finish, we multiply the number of valid $k$ by the number of valid $c$ to get $4 \cdot 21= \boxed{\textbf{(D) } 84}$

 \textbf{Note} You can also find the coefficients of the polynomial using Vietas Formulas. Because $f$ is monic, the leading coefficient of this quadratic is 1. The sum of the values is $6$, and the product of the values is $(3-\sqrt{5})(3+\sqrt{5})=3^2-(\sqrt{5})^2=9-5=4$. Thus by Vietas, resulting in the same polynomial, $x^2-6x+4$, but just a little bit quicker. ~Voidling

 



\textbf{Solution 5 (Quick Calculus)}

Let the polynomial be $f(x)=x^3+ax^2+bx+c$. First, note that if a cubic polynomial $f(x)$ has the same value at $a-b$ and $a+b$, then the derivative of the polynomial is always equal to $-b^2$. We can then take the derivative of the polynomial and make it equal to $-5$ at $x=3$.

 
Taking the derivative yields $f'(x)=3x^2+2ax+b=-5$. Plugging in $x=3$ gives $f'(3)=3(3)^2+2a(3)+b=-5$ and this gives the equation $27+6a+b=-5$, and therefore, $b=-6a-32$. Plugging this value into the given inequality yields:

 
 \[-10\le -6a-32\le 10\implies 22\le -6a\le 42\implies -7\le a\le -\frac{11}{3}\]Therefore, $a\in\{-7, -6, -5, -4\}$. Notice that these values of $a$ also satisfy the inequality. Then, we multiply the number of valid $a$ by the number of valid $c$ and this gives $4 \cdot 21= \boxed{\textbf{(D) } 84}$ total polynomials.

 



\textbf{Solution 6 (No Cubing Radicals)}

Let $f(x) = x^3 + ax^2 + bx + c$, and let $p = 3 + \sqrt{5}$ and $q = 3 - \sqrt{5}$. Note that $p + q = 6$ and $pq = 9 - 5 = 4$.

 Substituting $p$ and $q$ has the condition $f(p) = f(q)$ become
 \[p^3 + ap^2 + bp + c = q^3 + aq^2 + bq + c,\]so
 \[(p^3 - q^3) + a(p^2 - q^2) + b(p - q) = 0.\]

 Since $p \neq q$, we can divide by $p - q$ to get
 \[(p^2 + pq + q^2) + a(p + q) + b = 0.\]

 We have $p^2 + q^2 = (p + q)^2 - 2pq = 36 - 8 = 28$, so $p^2 + pq + q^2 = 28 + 4 = 32$. The equation becomes
 \[32 + 6a + b = 0, \quad\text{so}\quad b = -32 - 6a.\]

 We need $-10 \le b \le 10$, which means $-10 \le -32 - 6a \le 10$. This gives $-7 \le a \le -\frac{11}{3}$, so $a \in \{-7, -6, -5, -4\}$, with corresponding $b = 10, 4, -2, -8$, and note that each pair is unique.

 Finally, since the constant $c$ cancels out, $c$ can be any of the $21$ integers from $-10$ to $10$. The number of polynomials is $4 \cdot 21 = \boxed{\textbf{(D)}~84}$.

 


 \textbf{Note} If we solve for $a$ in terms of $b$ instead of solving for $b$ in terms of $a$ for the final inequality, we will get $-10 \le \frac{-32 - b}{6} \le 10$. Simplifying gives $-92 \le b \le 28$. This is entirely correct, however it is not useful in our situation given that $-10 \le b \le 10$. Aside from solving for $a$ in terms of $b$ being simpler, we also note that if we solve for $a$ in terms of $b$ and treat $a$ as $y$ and $b$ as $x$ in a graph, the slope will be $-\frac{1}{6}$. The boundaries the problem gives us are $-10 \le a \le 10$, so our slope that is close to $0$ will make the line travel for a long distance on the $b$ axis before reaching the other boundary, making it so that inequality will not provide much new information to us. If we do this quick check, we can realize that solving for $b$ in terms of $a$ is the correct thing to do.

 \textbf{Another Note} Due to $a$ and $b$ being integers, once we found the values of $a$ that work it is guaranteed that each value of $a$ has a corresponding value of $b$ which is just $-32-6a$, which is what we solved for, and that is an integer. We also know that $-32-6a$ is between $-10$ and $10$ because by definition it is the same as $b$ which has the same restriction.

 


	\item (\textit{AMC 12 12A 2026}) Six dogs find a large cookie in the shape of a regular hexagon $ABCDEF$. To divide the cookie, the dogs agree to choose point $P$ uniformly at random in the interior of the hexagon. The dogs will then cut out the segments $\overline{PA}$, $\overline{PB}$, $\overline{PC}$, $\overline{PD}$, $\overline{PE}$, and $\overline{PF}$, dividing the cooking into six triangles, one for each dog. A dog will be happy if and only if its piece has at least $\frac{1}{10}$ of the total area of the hexagon. What is the probability that all six dogs will be happy?


 $\textbf{(A)}\   \dfrac{\sqrt{3}}{50}  \qquad 	\textbf{(B)}\   \dfrac{1}{25}  \qquad 	\textbf{(C)}\   \dfrac{3\sqrt{3}}{50}  \qquad 	\textbf{(D)}\   \dfrac{4}{25}  \qquad 	\textbf{(E)}\   \dfrac{2}{5}$

 \textbf{Answer: }D



\textbf{Solution 1}

Let $O$ be the center and $S$ the total area. Write $[XYZ]$ for the area of $\triangle XYZ$. Extend $OP$ to meet the boundary at $Q$, and set $t=\frac{OP}{OQ}$.

 With a fixed base, a triangle's area changes linearly as its vertex moves along a line. Thus, for every side $UV$,
 \[[PUV]=(1-t)[OUV]+t[QUV] =(1-t)\frac S6+t[QUV] \ge(1-t)\frac S6.\]Equality holds for the side containing $Q$, so the smallest piece has area exactly $(1-t)S/6$. Therefore, all six dogs are happy precisely when
 \[\frac{(1-t)S}{6}\ge\frac S{10} \quad\Longleftrightarrow\quad t\le\frac25.\]Along every ray from $O$, the allowed points occupy the first $\frac{2}{5}$ of the distance to the boundary. Hence the allowed region is the blue hexagon below, a copy of the original scaled by $\frac{2}{5}$ about $O$.

 
                    \begin{figure}[H]
                        \centering
                        \includegraphics[width=0.5\textwidth]{Solution/AMC_12_12A_2026_21_Solution_Figure_1.png}
                    \end{figure}
                    

 Since $P$ is chosen uniformly, the probability equals the area ratio:
 \[\left(\frac25\right)^2=\boxed{\frac4{25}},\]We choose $\boxed{\textbf{(D)}}$.

 



\textbf{Solution 2}

We may assume WLOG that the hexagon has side length $1$. The area of the hexagon is then
 \[\frac{\sqrt{3}}{4}\cdot 6=\frac{3\sqrt{3}}{2}\]so the area that each dog must get in order to be happy is $\frac{3\sqrt{3}}{20}$. This implies that the distance $h$ from $P$ to each side, or the height of the triangle of each dog, must satisfy
 \[\frac{1}{2}\cdot 1\cdot h\ge \frac{3\sqrt{3}}{20}\implies h\ge \frac{3\sqrt{3}}{10}\]By symmetry, the area produced that yields a valid configuration is a smaller hexagon. The distance between opposite sides of the original hexagon can be found to be $\sqrt{3}$, so the distance between opposite sides of the new hexagon is
 \[\sqrt{3}-2\cdot \frac{3\sqrt{3}}{10}=\frac{2\sqrt{3}}{5}\]Therefore, the ratio between side lengths of the smaller hexagon to the larger hexagon is
 \[\frac{\frac{2\sqrt{3}}{5}}{\sqrt{3}}=\frac{2}{5}\]since the ratio of the distances between opposite sides is the same as the ratio of the lengths of the sides. Therefore, the ratio of the areas, which is precisely the fraction of the original hexagon that is valid, or the answer, is $\left(\frac{2}{5}\right)^2=\frac{4}{25}$, which is $\boxed{\textbf{(D)}}$.

 


 
Other Path you can take: Alternatively, instead of using the total distance between opposite sides, we can analyze the region from the center using the apothem. 

 Let the distances from the interior point $P$ to the sides of the hexagon be $h_1, h_2, \dots, h_6$. Since each base has a length of $1$, the total area of the hexagon is the sum of the areas of these six triangles:
 \[\sum_{k=1}^6 \frac{1}{2}h_k = \frac{3\sqrt{3}}{2} \implies \sum_{k=1}^6 h_k = 3\sqrt{3}\]As established above, each height must satisfy $h_k \ge \frac{3\sqrt{3}}{10}$. This condition bounds $P$ within a smaller, concentric regular hexagon shifted inward by a distance of $\frac{3\sqrt{3}}{10}$ from each side. 

 Since the apothem of the original unit hexagon is $a = \frac{\sqrt{3}}{2}$, the apothem of this smaller valid region is:
 \[a_{\text{small}} = a - \frac{3\sqrt{3}}{10} = \frac{\sqrt{3}}{2} - \frac{3\sqrt{3}}{10} = \frac{\sqrt{3}}{5}\]Using the ratio of their apothems, the ratio of their areas (and thus the probability) is:
 \[\text{Probability} = \left( \frac{a_{\text{small}}}{a} \right)^2 = \left( \frac{\frac{\sqrt{3}}{5}}{\frac{\sqrt{3}}{2}} \right)^2 = \left( \frac{2}{5} \right)^2 = \frac{4}{25}\]



	\item (\textit{AMC 12 12A 2026}) A flock of $8$ birds is flying in a formation, led by a single leader bird, which has exactly $2$ followers immediately behind. Each non-leader bird follows exactly one bird, and has $0$, $1$, or $2$ followers immediately behind. Two possible formations are shown in the picture below. (When a bird has two followers, swapping its left and right subflocks can produce a different formation, as shown in the figure.)


 

                    \begin{figure}[H]
                        \centering
                        \includegraphics[width=0.5\textwidth]{Problem/AMC_12_12A_2026_22_Problem_Figure_1.png}
                    \end{figure}
                    


 


 What is the total number of possible formations?


 $\textbf{(A)}\   48  \qquad 	\textbf{(B)}\   52  \qquad 	\textbf{(C)}\   76  \qquad 	\textbf{(D)}\   132  \qquad 	\textbf{(E)}\   196$

 \textbf{Answer: }C



\textbf{Solution 1}

We temporarily remove every bird with exactly one follower, connecting its predecessor directly to its follower. This preserves the left-right order of every branch.

 If $k$ birds remain with two followers, the reduced formation has $2k$ links and hence $2k+1$ birds. Thus $k=1,2,3$. There are respectively $1,2,5$ reduced shapes, shown below.

 For $k=3$, the other two branching birds lie on opposite sides of the leader in one way, or on the same side in $2\cdot2$ ways.

 
                    \begin{figure}[H]
                        \centering
                        \includegraphics[width=0.5\textwidth]{Solution/AMC_12_12A_2026_22_Solution_Figure_1.png}
                    \end{figure}
                    

 For each reduced shape, we reinsert the $7-2k$ removed birds along its $2k$ links. Choosing how many to place on each link is a stars-and-bars problem, giving
 \[\binom{(7-2k)+(2k)-1}{2k-1}=\binom{6}{2k-1}\]possibilities. No birds may be inserted above the leader, which must still have two immediate followers. Each original formation is recovered uniquely.

 Therefore the total is
 \[1\binom61+2\binom63+5\binom65 =6+40+30 =\boxed{\text{(C)}\ 76}.\]




\textbf{Solution 2}

Let $a_n$ denote the number of arrangements with $n$ birds, except that the condition that the leader is followed by two followers is removed. We perform recursion; we have two cases:

 
            \begin{itemize}
        \item When the leader is followed by one bird, then the follower and the rest of the birds form the case $a_{n-1}$.
        \item When the leader is followed by two birds, then say that the set of the left follower and all birds “behind” the left bird (i.e. birds that, when following the path formed by the bird it follows, eventually leads to the left bird) has size $k$. Then, the left bird and the birds behind it form $a_k$, and the right bird and the birds behind it form $a_{n-k-1}$. We know that $k$ ranges from $1$ to $n-2$ since the leader and right bird are not included in the set of birds on the left.
            \end{itemize}
            We therefore conclude that
 \[a_n=a_{n-1}+(a_{n-2}a_1+a_{n-3}a_2+a_{n-4}a_3+\cdots +a_1a_{n-2})\]Now we may apply our recursion. It is easy to find that $a_1=1$, $a_2=1$, $a_3=2$. Then,
 \[a_4=a_3+a_2a_1+a_1a_2=2+1\cdot 1+1\cdot 1=4\]
 \[a_5=a_4+a_3a_1+a_2a_2+a_1a_3=4+5=9\]
 \[a_6=a_5+a_4a_1+a_3a_2+a_2a_3+a_1a_4=9+12=21\]Finally, our answer for $n=8$ does not involve the first case above since the leader must be followed by two followers, so we subtract $a_7$ from $a_8$, which yields our answer:
 \[a_8-a_7=a_6a_1+a_5a_2+a_4a_3+a_3a_4+a_2a_5+a_1a_6=76\implies \boxed{\textbf{(C)}}\]




\textbf{Solution 3}

Let's solve this with recursion. Let $f(n)$ be the possible of formations with $n$ birds, where every bird can have $0$, $1$, or $2$ followers. 

 For $n\ge 2$, there are two cases: the leader has one follower or two followers.

 If the leader has one follower, then the remaining $n-1$ birds can form any valid formation, so there are $f(n-1)$ possibilities.

 If the leader has two followers, then if the left subflock has $k$ birds, the right subflock has $n-1-k$ birds. The left and right subflocks are distinguishable, so there are $f(k)f(n-1-k)$ possibilities. We will sum this for all possible valid $k$ values. 

 Our recurrence is 
 \[f(n)=f(n-1)+\sum_{k-1}^{n-2} f(k)f(n-1-k).\]

 Following this recurrence, we get the following table until $n=6$:
 \[\begin{array}{c|l|c} n & f(n)\\ \hline 1 & 1\\ 2 & 1\\ 3 & 2\\ 4 & 4\\ 5 & 9\\ 6 & 21 \end{array}\]

 We don't continue beyond $n=6$ because $n=8$ is different, since the leader must have exactly $2$ followers. Therefore, we must exclude the case where the leader has exactly $1$ follower. This means we need:

 
 \[\sum_{k=1}^6 f(k)f(7-k) = (1)(21)+(1)(9)+(2)(4)+(4)(2)+(9)(1)+(21)(1) = 76.\]

 Therefore, our answer is $\boxed{76}$.

 


 



\textbf{Solution 4 (Lattice Path)}

Instead of counting the formations directly, we transform each formation into a lattice path.

 Let $a_n$ denote the number of unfilled positions after placing $n$ birds. We fill these positions one at a time, completing each left subflock before its corresponding right subflock. If a bird has only one follower, no left-right distinction is made.

 Since the leader has exactly two followers, we have
 \[a_1=2.\]

 Each subsequent bird occupies one unfilled position and creates either $0$, $1$, or $2$ new positions. Thus,
 \[a_{n+1}= \begin{cases} a_n-1 & \text{if the bird has no followers},\\ a_n & \text{if the bird has one follower},\\ a_n+1 & \text{if the bird has two followers}. \end{cases}\]

 Consequently, each formation corresponds to a lattice path with steps $(1,-1)$, $(1,0)$, and $(1,1)$. Conversely, each such valid path uniquely determines a formation, since the number of followers of each successive bird is determined by the change in $a_n$, and the order in which positions are filled is fixed.

 Since the formation cannot be completed before all eight birds are placed, we require
 \[a_n>0\quad\text{for }1\leq n<8,\qquad a_8=0.\]

 Furthermore, since each step can increase or decrease $a_n$ by at most $1$, the initial and final conditions imply
 \[a_n\leq n+1,\qquad a_n\leq 8-n.\]

 Hence, all intermediate lattice points must lie in the region
 \[1\leq a_n\leq\min(n+1,8-n).\]

 The shaded region represents the admissible intermediate states. Each node is labeled with the number of valid paths reaching it.

 
                    \begin{figure}[H]
                        \centering
                        \includegraphics[width=0.5\textwidth]{Solution/AMC_12_12A_2026_22_Solution_Figure_2.png}
                    \end{figure}
                    

 

 The pink node represents the initial state, while the green nodes indicate the final count of valid paths. In particular, the terminal point $(8,0)$ is labeled $76$. Therefore, the total number of possible formations is
 \[\boxed{\textbf{(C)}\ 76}.\]



	\item (\textit{AMC 12 12A 2026}) There is a unique real number $x$ between $0$ and $1$ such that $x^{3} + x^{2} = 1$. For that $x$, there exist unique integers $a$, $b$, and $c$ such that the infinite product

 \[(x+1)(x^{2}+1)(x^{4}+1)(x^{8}+1)(x^{16}+1) \cdots\]
equals $ax^{2} + bx + c$. What is $a+b+c$?


 $\textbf{(A)}\   2  \qquad 	\textbf{(B)}\   3  \qquad 	\textbf{(C)}\   4  \qquad 	\textbf{(D)}\   5  \qquad 	\textbf{(E)}\   6$

 \textbf{Answer: }D



\textbf{Solution}

The infinite product can be represented as
 \[\lim_{n\to\infty}\prod_{k=0}^n (1+x^{2^k})\]When we multiply the expression by $(1-x)$, the entire expression telescopes, yielding
 \[\lim_{n\to\infty}\frac{1-x^{2^{n+1}}}{1-x}\]which goes to $\frac{1}{1-x}$ as $n\to\infty$.

 Finally, we equate to get $1=(1-x)(ax^2+bx+c)$. Expanding yields
 \[ax^3+x^2(b-a)+x(c-b)+(1-c)=0\]We can now guess values for $a$. Coincidentally, $a=1$ works, yielding $b=2$ and $c=2$. Hence the answer is $1+2+2=5$, or $\boxed{\textbf{(D)}}$.

 



\textbf{Solution 1.5}

As in solution 1, we get $ax^3+x^2(b-a)+x(c-b)=c-1$. Instead of guessing, we can match this equation to $x^3+x^2=1$ by ensuring the coefficients are the same. So, we get $a=1$, $b-a=1$, $c-b=0$, and $c-1=1$. Thus, $a=1$, $b=2$, and $c=2$, so the answer is $1+2+2=5$, or $\boxed{\textbf{(D)}}$.

 Technically you can’t do this because it is not an identity, but because x is irrational, you can kind of get away with it- this is still mathematically wrong. Solution 2 is a more complete way.- Skip\_God (not my solution)

 



\textbf{Solution 2}

Continuing as in Solution 1 and getting the polynomial
 \[ax^3 + (b-a)x^2 + (c-b)x + (1-c),\]we add and subtract $ax^2$ to obtain
 \begin{align*} ax^3 + ax^2 - ax^2 + (b-a)x^2 + (c-b)x + (1-c) &= a(x^3 + x^2) + (b-2a)x^2 + (c-b)x + (1-c) \\ &= (b-2a)x^2 + (c-b)x + (1+a-c). \end{align*}
We assume that $x$ and $x^2$ are both irrational, so the only way for this polynomial to equal $0$ is for each of the coefficients and the constant term to be $0$. This sets up the system of equations

 
 \begin{align*} b-2a &= 0\\ c-b &= 0\\ 1+a-c &= 0, \end{align*}
whence $a=1$, $b=2$, and $c=2$, so $a + b + c = \boxed{\textbf{(D) }5}$

 Note: No need to assume that $x^2$ and $x$ are both irrational; we can prove it by using the handy Rational Root Theorem. It states that since the polynomial $x^3+x^2-1=0$ has all rational coefficients, any rational root (of the form $\frac p q$) of the polynomial must have $p | (-1)$ (-1 is from the last term of the polynomial) and $q | 1$ (1 is from the leading coefficient of the polynomial). We try out the only roots that satisfy these conditions, $\pm1$, and they don't work. Therefore, $x$ and $x^2$ are both irrational.~ WhatDoHumanitariansEat

 



\textbf{Solution 3}

We rewrite the given equation $x^3 + x^2 = 1$ into these two equations:
 \begin{align}  x^n + x^{n-1} = x^{n-3} \implies x^n = x^{n-3} - x^{n-1}, n \in \mathbb{R} \\ x^{m+1} + x^m = x^{m-2} \implies x^m = x^{m-2} - x^{m+1}, m \in \mathbb{R} \end{align}Note that the infinite product can be written as $\sum_{n=0}^\infty x^n = 1 + x + x^2 + x^3 + x^4 + x^5...$

 If we replace every $x^{4k}$ term with $x^{4k-2}-x^{4k+1}, k \in \mathbb{R}$ (equation 1) and

 we replace every $x^{4k+2}$ term with $x^{4k} - x^{4k+3}, k \in \mathbb{R}$ (equation 2),

 every section of

 
 \begin{align*} \\ x^{4k}+x^{4k+1}+x^{4k+2}+x^{4k+3} \\ = x^{4k-2} - x^{4k+1} + x^{4k+1} + x^{4k} - x^{4k+3} + x^{4k+3} \\ = x^{4k-2} + x^{4k} \\ \end{align*}


 therefore:

 
 \begin{align*}\\ \sum_{n=0}^\infty x^n = 1 + x + \sum_{m=0}^\infty x^{2m}\\ \end{align*}


 Let:
 \begin{align*}\\ u = \sum_{m=0}^\infty x^{2m}\\ \end{align*}


 So:
 \begin{align*}\\ \sum_{n=0}^\infty x^n = \sum_{n=0}^\infty x^{2n} + \sum_{n=0}^\infty x^{2n+1} = \sum_{n=0}^\infty x^{2n} + x \sum_{n=0}^\infty x^{2n} = u + ux\\ \end{align*}


 And now solving for $u$:
 \begin{align*}\\ u + ux = 1 + x + u \implies ux = 1+x \implies u = \frac{1}{x} + 1\\ \end{align*}


 Therefore:
 \begin{align*}\\ u + ux = \frac{1}{x} + 1 + (x)(\frac{1}{x} + 1) = \frac{1}{x} + 1 + 1 + x\\ \end{align*}


 Because $1 = x^3 + x^2$, $\frac{1}{x} = x^2 + x$

 Therefore:
 \begin{align*}\\ \frac{1}{x} + 1 + 1 + x = x^2 + x + 2 + x = x^2 + 2x + 2\\ \end{align*}


 And $1 + 2 + 2 = \boxed{\textbf{(D) }5}$




\textbf{Solution 4}

Note that the provided product
 \[(x + 1)(x^2 + 1)(x^4 + 1)(x^8 + 1)(x^{16} + 1)\dotsm\]expands as a formal power series
 \[1 + x + x^2 + x^3 + x^4 + x^5 + \dotsb,\]which converges to ${\frac{1}{1 - x}}$ for the given $x \in (0, 1)$. We wish to find a polynomial $ax^2 + bx + c$ such that $(1 - x)(ax^2 + bx + c) = 1$.

 Recall we are given $x^3 + x^2 = 1$, which can be rearranged as
 \[x^3 = 1 - x^2 = (1 - x)(1 + x), \qquad \text{and} \qquad x^2 = 1 - x^3 = (1 - x)(1 + x + x^2).\]Adding,
 \[x^3 + x^2 = (1 - x)((1 + x) + (1 + x + x^2)) = (1 - x)(2 + 2x + x^2),\]which is $1$ since $x^3 + x^2 = 1$.

 We are done, having found $a = 1$, $b = 2$, $c = 2$, so $a + b + c = \boxed{\textbf{(D) }5}$.

 



\textbf{Solution 5}

As Solution 4 showed, 
 \[(x + 1)(x^2 + 1)(x^4 + 1)(x^8 + 1)(x^{16} + 1)\dotsm\]is rearranged as 
 \[1 + x + x^2 + x^3 + x^4 + x^5 + \dotsb,\]and it is converged to $\frac{1}{1 - x}$.

 Using the relation $x^3 + x^2 = 1$, we can express:
 \[\frac{x^3 + x^2}{1-x} = \frac{1}{1-x}\]

 By performing polynomial long division on $\frac{x^3+x^2}{1-x}$:
 \[\frac{x^3+x^2}{-(x-1)} = -x^2 - 2x - 2 + \frac{2}{1-x}\]

 Equating this to $\frac{1}{1-x}$:
 \[\frac{1}{1-x} = -x^2 - 2x - 2 + \frac{2}{1-x}\]

 Rearranging the terms yields:
 \[\frac{1}{1-x} = x^2 + 2x + 2\]

 Thus, $(x + 1)(x^2 + 1)(x^4 + 1)(x^8 + 1)(x^{16} + 1)\dotsm$ equals $x^2 + 2x + 2$.

 And $1 + 2 + 2 = \boxed{\textbf{(D) }5}$.

 



\textbf{Solution 6}

As used in Solutions 4 and 5, we can set $S$ as the sum of our infinite series
 \[S = 1 + x + x^2 + x^3 + x^4 + x^5 + \dotsb\]Using the relation we are given, we can extrapolate the following infinite list of equations:
 \[x^3 + x^2 = 1\]
 \[x^4 + x^3 = x\]
 \[x^5 + x^4 = x^2\]
 \[\vdots\]Summing them, we get
 \[x^2 + 2(x^3 + x^4 + x^5 + \dotsm) = 1 + x + x^2 + \dotsm\] which can be rearranged and substituted to get
 \[x^3 + x^4 + x^5 + \dotsm =  \dfrac{S - x^2}{2}\]Our original series can now be written as 
 \[S = 1 + x + x^2 + \dfrac{S - x^2}{2}\]Solving for $S$ gives us 
 \[S = x^2 + 2x + 2\]And $1 + 2 + 2 = \boxed{\textbf{(D) }5}$




\textbf{Solution 7}

As shown in solution 4,
 \[(x + 1)(x^2 + 1)(x^4 + 1)\dotsm\]converges to ${\frac{1}{1 - x}}$
$x^3 + x^2 = 1$ can be rearranged as 
 \[x^3 = 1 - x^2\] which can be factored as 
 \[x^3 = (1 - x)(1 + x)\] and rearranged as 
 \[1 - x = \frac{x^3}{1+x}\]Substituting $1 - x$ in the denominator with $\frac{x^3}{1+x}$ yields 
 \[\frac{1 + x}{x^3} = ax^2 + bx + c\]Since $1 = x^3 + x^2$, 
 \[1 + x = (1 + x)(x^3 + x^2) = x^4 + 2x^3 + x^2(x^3 + x^2) = x^5 + 2x^4 + 2x^3\]

 Dividing by $x^3$ yields 
 \[x^2 + 2x + 2\] And $1 + 2 + 2 = \boxed{\textbf{(D) }5}$



	\item (\textit{AMC 12 12A 2026}) Trapezoid $ABCD$ has $AB = 4$, $CD = 5$, and $\angle DAB = \angle ABC = 90^{\circ}$. Circles $\omega_{1}$ and $\omega_{2}$ lie inside the trapezoid as follows: $\omega_{1}$ is tangent to $\overline{AB}$, $\overline{AD}$, and $\overline{DC}$; $\omega_{2}$ is tangent to $\overline{AB}$, $\overline{BC}$, and $\overline{DC}$; and $\omega_{1}$ and $\omega_{2}$ are tangent to each other. The distance between the centers of $\omega_{1}$ and $\omega_{2}$ is $m - n\sqrt{p}$, where $m$, $n$, and $p$ are positive integers such that $p$ is not divisible by the square of any prime. What is $m + n + p$?


 $\textbf{(A)}\   30  \qquad 	\textbf{(B)}\   47  \qquad 	\textbf{(C)}\   53  \qquad 	\textbf{(D)}\   62  \qquad 	\textbf{(E)}\   74$

 \textbf{Answer: }D



\textbf{Solution 1}

Let the circles have centers $O_1,O_2$ and radii $r_1,r_2$. We write $d=O_1O_2$. Since the circles are externally tangent,
 \[d=r_1+r_2.\]

 Let $P,Q$ be the feet of the perpendiculars from $O_1,O_2$ to $AB$. Tangency to $AD$ and $BC$ gives $AP=r_1$ and $QB=r_2$. Thus
 \[PQ=AB-AP-QB=4-r_1-r_2=4-d.\]

 Extend $AB$ and $CD$ to meet at $E$, and let $\theta=\angle DEB$. Each center is equally distant from the lines $AB$ and $CD$, so both centers lie on the same angle bisector. Therefore, $O_1O_2$ makes an angle of $\frac{\theta}{2}$ with $AB$.

 
                    \begin{figure}[H]
                        \centering
                        \includegraphics[width=0.5\textwidth]{Solution/AMC_12_12A_2026_24_Solution_Figure_1.png}
                    \end{figure}
                    

 Since $AD$ and $BC$ are perpendicular to $AB$, the projection of $CD$ onto $AB$ has length $4$. Hence
 \[\cos\theta=\frac{4}{5}, \qquad \cos\frac{\theta}{2} =\sqrt{\frac{1+\frac45}{2}} =\frac3{\sqrt{10}}.\]

 Projecting $O_1O_2$ onto $AB$ gives
 \[4-d=PQ=d\cos\frac{\theta}{2} =\frac{3d}{\sqrt{10}}.\]Solving,
 \[d=\frac{4\sqrt{10}}{\sqrt{10}+3} =40-12\sqrt{10}.\]Thus $m+n+p=40+12+10=\boxed{62}$, we choose $\boxed{\textbf{(D)}}$.

 



\textbf{Solution 2 (Constructions)}

\textbf{Diagram} 
                    \begin{figure}[H]
                        \centering
                        \includegraphics[width=0.5\textwidth]{Solution/AMC_12_12A_2026_24_Solution_Figure_2.png}
                    \end{figure}
                    




\textbf{Solution}

Let $r_1$ and $r_2$ be the radii of $\omega_1$ and $\omega_2$, with centers $O_1$ and $O_2$. Extend $BA$ and $CD$ to meet at $E$. Since both circles are tangent to $EB$ and $EC$, both centers lie on the angle bisector of $\angle BEC$. Let this bisector meet $BC$ at $K$, and let $\angle BEK = \angle KEC = \theta$.

 Now, drop the perpendicular from $O_1$ to $BC$, and draw the line through $O_2$ parallel to $BC$. Let these lines meet at $F$. Since the circles are externally tangent at $T$, and both are tangent to $AB$,
 \[O_1O_2 = r_1 + r_2, \qquad O_2F = r_2 - r_1, \qquad O_1F = 4 - r_1 - r_2.\]

 Let $P$ be the projection of $D$ onto $BC$. Then $DP = AB = 4$ and $DC = 5$, so $\triangle DPC$ is a $3$-$4$-$5$ triangle with $PC = 3$. Since $DP \parallel EB$, we have $\triangle EBC \sim \triangle DPC$, so $EB : BC : EC = 4 : 3 : 5$.

 By the Angle Bisector Theorem on $\triangle EBC$,
 \[\frac{BK}{KC} = \frac{EB}{EC} = \frac{4}{5}, \quad\text{so}\quad BK = \frac{4}{9}BC.\]Then $\frac{BK}{EB} = \frac{4}{9}\cdot\frac{3}{4} = \frac{1}{3}$, so $\triangle EBK$ is a $1 : 3 : \sqrt{10}$ right triangle with $BK : EB : EK = 1 : 3 : \sqrt{10}$.

 Since $O_1F \parallel EB$, we have $\angle FO_1O_2 = \angle BEK = \theta$, so $\triangle O_1FO_2 \sim \triangle EBK$. Therefore
 \[O_2F : O_1F : O_1O_2 = 1 : 3 : \sqrt{10}.\]

 Let $d = r_1 + r_2$. Using $O_1F : O_1O_2 = 3 : \sqrt{10}$,
 \[\frac{4 - d}{d} = \frac{3}{\sqrt{10}} \implies 4\sqrt{10} = d\left(\sqrt{10} + 3\right) \implies d = \frac{4\sqrt{10}}{\sqrt{10} + 3}.\]Rationalizing the denominator,
 \[d = 4\sqrt{10}\left(\sqrt{10} - 3\right) = 40 - 12\sqrt{10}.\]

 Thus $O_1O_2 = 40 - 12\sqrt{10}$, so $m + n + p = 40 + 12 + 10 = \boxed{\textbf{(D)}~62}$.

 \textbf{Remark} This was one of the coolest geo problems I've seen on the AMCs in a good while. The solution above was my solution during the actual test. Eww coordinates :D

 



\textbf{Solution 3}

Let $A=(0,0)$, $B=(4,0)$, $D=(0,h)$, and $C=(4,h+3)$.

 Since $AB=4$ and $CD=5$, the Pythagorean theorem gives the vertical difference between $C$ and $D$ as $3$.

 The equation of line $DC$ is
 \[3x-4y+4h=0.\]

 Let the radii of $\omega_1$ and $\omega_2$ be $r_1$ and $r_2$, respectively.

 Since $\omega_1$ is tangent to both $AB$ and $AD$, its center is $(r_1,r_1)$.

 Since $\omega_2$ is tangent to both $AB$ and $BC$, its center is $(4-r_2,r_2)$.

 Using the distance formula for a point and a line, the tangency of $\omega_1$ to $DC$ gives

 
 \[\frac{3r_1-4r_1+4h}{5}=r_1.\]

 Simplifying,

 
 \[4h-r_1=5r_1\]

 \[r_1=\frac{2h}{3}.\]

 Similarly, the tangency of $\omega_2$ to $DC$ gives

 
 \[\frac{3(4-r_2)-4r_2+4h}{5}=r_2.\]

 Simplifying,

 
 \[12+4h=12r_2\]

 \[r_2=1+\frac{h}{3}.\]

 Since the circles are tangent to each other, the distance between their centers equals the sum of their radii. Thus,

 
 \[(4-r_1-r_2)^2+(r_2-r_1)^2 =(r_1+r_2)^2.\]

 Expanding and simplifying,

 
 \[(r_2-r_1)^2=8(r_1+r_2)-16.\]

 Substituting our expressions for the radii,

 
 \[\left(1-\frac{h}{3}\right)^2 =8(h+1)-16.\]

 Multiplying by $9$ and simplifying,

 
 \[h^2-78h+81=0.\]

 Using the quadratic formula,

 
 \[h=\frac{78\pm\sqrt{78^2-4(81)}}{2}\]

 \[h=39\pm12\sqrt{10}.\]

 Since the circles must fit inside the trapezoid, we take the smaller root:

 
 \[h=39-12\sqrt{10}.\]

 Finally, the distance between the centers is

 
 \[r_1+r_2=1+h.\]

 Therefore,

 
 \[d=40-12\sqrt{10}.\]

 Hence, $m=40$, $n=12$, and $p=10$, so

 
 \[m+n+p=40+12+10=\boxed{62}.\]

 The answer is $\boxed{\text{(D) }62}$.$\sim$Cyrus825

 


	\item (\textit{AMC 12 12A 2026}) Let $m$ be a positive integer such that there are at least $50$ positive integers $x$ less than $m$ for which $x^2 - 1$ is a multiple of $m$. What is the sum of the digits of the least possible value of $m$?


 $\textbf{(A)}\   12  \qquad 	\textbf{(B)}\   15  \qquad 	\textbf{(C)}\   16  \qquad 	\textbf{(D)}\   18  \qquad 	\textbf{(E)}\   22$

 \textbf{Answer: }B



\textbf{Solution 1}

Let $p^k$ be the largest power of a prime $p$ that divides $x^2-1$. We focus on the exponent of each prime involved. First, if $p\ge 3$, then we have
 \[x^2\equiv 1\pmod{p^k}\implies (x+1)(x-1)\equiv 0\pmod{p^k}\]But $x+1$ and $x-1$ cannot both be divisible by $p$ since then $x+1\equiv x-1\equiv 0\pmod{p}\implies 2\equiv 0\pmod{p}$, which is false. Therefore, all powers of $p$ must be “stored” in either $x+1$ or $x-1$, hence $x\equiv \pm 1\pmod{p^k}$. As a result, each odd prime power contributes a multiplication of $2$, or a doubling, to the total number of solutions by the Chinese Remainder Theorem.

 For $p=2$, we may use the above argument, except that both $x+1$ and $x-1$ may be divisible by $2$, but only one may be divisible by $4$. Therefore, one of $x+1$ and $x-1$ is even, and the other is divisible by $2^{k-1}$. This implies that $x\equiv \pm 1\pmod{2^{k-1}}$, so at most $4$ solutions exist to $x^2\equiv 1\pmod{2^k}$. We perform casework on the precise value of $k$.

 When $k=0$, there are no solutions created (clear). When $k=1$, then there is exactly one solution contribution, so the number of solutions does not increase. When $k=2$, then $x\equiv \pm 1\pmod{2}$, hence $x\equiv 1,3\pmod{4}$, resulting in a doubling of the number of solutions. Finally, when $k\ge 3$, then $x\equiv 1,2^{k-1}-1,2^{k-1}+1,2^k-1\pmod{2^k}$, resulting in a quadrupling of the number of solutions.

 From our analysis above, we see that $m$ being divisible by two copies of the same odd prime does not increase the number of solutions, hence $m$ should be squarefree with respect to odd primes. We also know that we should have $k\le 3$ for $p=2$, otherwise we are wasting powers of $2$; we also do not have $k=1$ since $k=0$ is better. Finally, we perform casework on the value of $k$ for $p=2$ (also known as $\nu_2$):

 
            \begin{itemize}
        \item When $k=0$, the minimum is $3\cdot 5\cdot 7\cdot 11\cdot 13\cdot 17$, which produces $2^6=64$ solutions.
        \item When $k=2$, the minimum is $2^2\cdot 3\cdot 5\cdot 7\cdot 11\cdot 13$, which produces $2^6=64$ solutions.
        \item When $k=3$, the minimum is $2^3\cdot 3\cdot 5\cdot 7\cdot 11$, which produces $4\cdot 2^4=64$ solutions.
            \end{itemize}
            By direct comparison, we see that $k=3$ produces the smallest $m$. Evaluating directly yields $9240$. The sum of the digits is $15$, or $\boxed{\textbf{(B)}}$.

 



\textbf{Solution 2}

We want the least $m$ such that at least 50 integers $x$, with $1\le x<m$, satisfy

 
 \[x^2\equiv 1\pmod m.\]

 Notice that $x^2-1=(x-1)(x+1)$. Since $x-1$ and $x+1$ differ by $2$, they cannot both be divisible by an odd prime.

 Thus, for every odd prime power dividing $m$, there are exactly $2$ possible residues: $x\equiv 1$ or $-1$.

 Let us now consider the powers of $2$. Looking at modulo $8$, we see that all four odd residues $1,3,5,7$ satisfy $x^2\equiv 1\pmod 8$.

 For higher powers $2^a$, notice that exactly one of $x-1$ and $x+1$ is divisible by $4$, while the other has exactly one factor of $2$. Thus, if $2^a\mid(x-1)(x+1)$, we must have

 
 \[x\equiv\pm1\pmod{2^{a-1}}.\]

 Each residue gives at most two possibilities modulo $2^a$, so there can never be more than $4$ solutions. 

 By the Chinese Remainder Theorem, if $m$ is divisible by $8$ and has $k$ distinct odd prime factors, then there are $4\cdot 2^k$ solutions modulo $m$. Indeed, for all $a \geq 3$, the four residues $1, -1, 1+2^{a-1}, -1+2^{a-1}$ all satisfy the congruence modulo $2^a$.

 We need at least 50 solutions, so we want

 
 \[4\cdot 2^k\ge 50.\]

 The smallest possible $k$ is $4$, giving $64$ solutions.

 To minimize $m$, we choose the four smallest odd primes, $3,5,7,11$, along with $8$. Thus,

 
 \[m=8\cdot3\cdot5\cdot7\cdot11=9240.\]

 The sum of the digits of $9240$ is $\boxed{15}.$




\textbf{Solution 3}

First, we try
 \[m_0=8\cdot3\cdot5\cdot7\cdot11=9240.\]Modulo $8$, all four odd residues work; modulo each odd prime, the two residues $\pm1$ work. The Chinese Remainder Theorem gives $4\cdot2^4=64$ solutions, so the least possible $m$ satisfies $m\le9240$.

 
                    \begin{figure}[H]
                        \centering
                        \includegraphics[width=0.5\textwidth]{Solution/AMC_12_12A_2026_25_Solution_Figure_1.png}
                    \end{figure}
                    

 To prove minimality, note that $\gcd(x-1,x+1)\mid2$. Thus each odd prime power $p^e$ allows exactly two roots, $x\equiv\pm1\pmod{p^e}$. A power $2^a$ allows at most four: for $a\ge3$, one of $x-1,x+1$ has exactly one factor of $2$, forcing
 \[x\equiv\pm1\pmod{2^{a-1}}.\]For $a\le2$, there are at most two roots.

 Let $k$ be the number of distinct odd prime factors of $m$. If $k\le3$, there are at most $4\cdot2^3=32$ solutions. If $k\ge5$, then
 \[m\ge3\cdot5\cdot7\cdot11\cdot13=15015>9240,\]also impossible. Hence $k=4$.

 Since $2\cdot2^4=32<50$, the power of $2$ must contribute more than two roots, requiring $8\mid m$. Therefore
 \[m\ge8\cdot3\cdot5\cdot7\cdot11=9240.\]Equality is attainable by our construction, so the answer is
 \[9+2+4+0=\boxed{\text{(B)}\ 15}.\]



\end{enumerate}